\documentclass{book}
\usepackage[margin=1in]{geometry}
\usepackage{hyperref}
\usepackage{ulem}
\usepackage{tipa} % Add this line to support \textipa command
\hypersetup{
    colorlinks=true,
    linkcolor=black,
    filecolor=magenta,      
    urlcolor=cyan,
    }

\urlstyle{same}

\title{French Conjugation}
\author{Junzhang Luo}

\begin{document}

\tableofcontents

\chapter{All Tenses}
The verb \textit{mange} conjugated with \textit{je} 
    \begin{enumerate}
        \item Present tenses:
            \begin{enumerate}
                \item Présent: je \textcolor{red}{mange} (I eat) / (I am eating)
                \item Présent Negation: je ne \textcolor{red}{mange} pas (I don't eat)
            \end{enumerate}
        
        \item Past tenses:
            \begin{enumerate}
                \item Passé Composé: J'\textcolor{red}{ai mangé} (I ate) / (I have eaten)
                \item Passé Composé Negation: Je n'\textcolor{red}{ai} pas  \textcolor{red}{mangé} (I didn't eat)
                \item COD - DOP: 
                    \begin{center}
                        J'\textcolor{red}{ai mangé} une pomme $\rightarrow$ Je l'\textcolor{red}{ai mangée} (I ate an apple $\rightarrow$ I ate it) \\
                        J'\textcolor{red}{ai mangé} deux pommes $\rightarrow$ Je les \textcolor{red}{ai mangées} \\
                        (I ate two apples $\rightarrow$ I ate them) \\
                        Je ne les \textcolor{red}{ai} pas \textcolor{red}{mangées}
                    \end{center}
                \item Imparfait: Je \textcolor{red}{mangeais} (I was eating) / (I used to eat)
                \item Plus-Que-Parfait: J'\textcolor{red}{avais mangé} (I had eaten)
                \item Passé Récent: Je \textcolor{red}{viens de manger} (I just ate)
                \item Passé Simple: Je \textcolor{red}{mangeai} (I ate)
            \end{enumerate}
        
        \item Future tenses:
            \begin{enumerate}
                \item Futur Simple: Je \textcolor{red}{mangerai} (I will eat)
                \item Futur Proche: Je \textcolor{red}{vais manger} (I am going to eat)
                \item Futur Antérieur: J'\textcolor{red}{aurai mangé} (I will have eaten)
            \end{enumerate}

        \item Impératif:
            \begin{center}
                \textcolor{red}{Mange} - Eat \\
                \textcolor{red}{Mange} une pomme \\
                \textcolor{red}{Mange}-la \\
                Ne la \textcolor{red}{mange} pas
            \end{center}

        \item Conditionnel:
            \begin{enumerate}
                \item Conditionnel Présent: Je \textcolor{red}{mangerais} (I would eat)
                \item Conditionnel Passé: J'\textcolor{red}{aurais mangé} (I would have eaten)
            \end{enumerate}

        \item Subjonctif:
            \begin{center}
                Il faut que je \textcolor{red}{mange} (I have to eat)
            \end{center}
    \end{enumerate}


\chapter{Past Tenses}

\section{La Phrase Simple}
For simple phrases, we have the following structure:
\begin{center}
    \textbf{Subject + conjugated verb + object + noun}
\end{center}
For example: Tu regardes un film. (You watch a movie)

\newpage

\section{Les Articles Indefinis}
The indefinite articles in French are:
\begin{itemize}
    \item Un (masculine singular)
    \item Une (feminine singular)
    \item Des (plural)
\end{itemize}

\newpage

\section{Passé Composé}

\subsection{Les Verbes Conjugés avec Avoir et Être}
Use Avoir when we have a direct object, and Être when we don't. For example:
\begin{itemize}
    \item Descendre (to go down):
    \begin{itemize}
        \item Être: Je \textcolor{red}{suis descendu} (I went downstairs)
        \item Avoir: J'\textcolor{red}{ai descendu} les valises (I took down the suitcases)
    \end{itemize}
    The difference here is we can have a direct object with Avoir but we can't with Être.

    \item Monter (to go up):
    \begin{itemize}
        \item Être: Il \textcolor{red}{suis monté} (He went upstairs)
        \item Avoir: Il \textcolor{red}{a monté} la valise au grenier (He took the suitcase to the attic)
    \end{itemize}
    The first one does not allow conjugation with a direct object, but the second one does.

    \item Rentrer (to go home):
    \begin{itemize}
        \item Être: Elle \textcolor{red}{est rentrée} à la maison (She went home)
        \item Avoir: Elle \textcolor{red}{a rentré} la voiture (She took the car home inside)
    \end{itemize}
    
    \item Sortir (to go out):
    \begin{itemize}
        \item Être: Vous \textcolor{red}{êtes sorti} dehor (You went outside)
        \item Avoir: Vous \textcolor{red}{avez sorti} le chien (You took the dog outside)
    \end{itemize}
\end{itemize}

\subsection{Definition}

The Passé Composé in French covers the following tenses in English:
\begin{enumerate}
    \item The Present Perfect: I have eaten
    \item The Simple Past: I ate
    \item The Emphatic Past: I did eat
\end{enumerate}
If we take the verb \textit{manger} (to eat) as an example, the Passé Composé is 
\begin{center}
\textbf{j'ai mangé} (I have eaten)
\end{center}

\subsection{Quand Utiliser le Passé Composé}
Passé Composé is all about actions started in the past and completed in the past. 

\begin{itemize}
    \item Something that took place in the past with a time reference. \\
    Some indicative words include:
    \begin{center}
        \begin{tabular}{|c|c|}
            \hline
            \textbf{French} & \textbf{English} \\
            \hline
            Hier & Yesterday \\
            \hline
            Il y a deux jours & Two days ago \\
            \hline
            L'été dernier & Last summer \\
            \hline
            Mercredi dernier & Last Wednesday \\
            \hline
            Ce matin & This morning \\
            \hline
        \end{tabular}
    \end{center}
    Examples:
    \begin{center}
        \begin{tabular}{|c|c|c|}
            \hline
            \textbf{Word} & \textbf{French} & \textbf{English} \\
            \hline
            Arriver (être) & Il \textcolor{red}{est arrivé} ce matin & He arrived this morning \\
            \hline
            Pleuvoir (avoir) & Il \textcolor{red}{a plu} la semaine dernière & It rained last week \\
            \hline
        \end{tabular}
    \end{center}
    \item To describe an event that usually only happens once.
    \begin{center}
        \begin{tabular}{|c|c|}
            \hline
            \textbf{French} & \textbf{English} \\
            \hline
            Mourir & to die \\
            \hline
            Naître & to be born \\
            \hline
            Se marier & to get married \\
            \hline
            Avorter & to abort \\
            \hline
            Avoir ... ans & to turn ... years old \\
            \hline
        \end{tabular}
    \end{center}
    Examples:
    \begin{center}
        \begin{tabular}{|c|c|c|}
            \hline
            \textbf{Word} & \textbf{French} & \textbf{English} \\
            \hline
            Mourir (être) & Il \textcolor{red}{est mort} l'année dernière & He died last year \\
            \hline
            Avoir (avoir) & J'\textcolor{red}{ai} récemment \textcolor{red}{eu} 30 ans & I recently turned 30 \\
            \hline
        \end{tabular}
    \end{center}

    \item To talk about something that happened repeatedly or at least a few times.
    \begin{center}
        \begin{tabular}{|c|c|}
            \hline
            \textbf{Word} & \textbf{French} \\
            \hline
            Visiter (avoir) & J'\textcolor{red}{ai visité} le Canada plusieurs fois \\
            \hline
        \end{tabular}
    \end{center}

    \item To indicate something that happened suddenly.\\
    Some indicative words include:
    \begin{center}
        \begin{tabular}{|c|c|}
            \hline
            \textbf{French} & \textbf{English} \\
            \hline
            Soudain & Suddenly \\
            \hline
            Tout à coup & All of a sudden \\
            \hline
            Tout de suite & Right away \\
            \hline
            Juste quand & Just when \\
            \hline
            D'un coup & Out of the blue \\
            \hline
        \end{tabular}
    \end{center}
    Examples:
    \begin{center}
        \begin{tabular}{|c|c|}
            \hline
            \textbf{Word} & \textbf{French} \\
            \hline
            Arriver (être) & La police \textcolor{red}{est arrivée} tout de suite \\ 
            & (The police arrived right away) \\
            \hline
            Disparaître (avoir) & Le magicien \textcolor{red}{a disparu} tout à coup \\ 
            & (The magician disappeared all of a sudden) \\
            \hline
        \end{tabular}
    \end{center}

    \item To talk about a series of action taking place successively. \\
    Some indicative words include:
    \begin{center}
        \begin{tabular}{|c|c|}
            \hline
            \textbf{French} & \textbf{English} \\
            \hline
            D'abord & First \\
            \hline
            Enfin & Finally \\
            \hline
            Ensuite & Then \\
            \hline
            Puis & Next \\
            \hline
            Après & After \\
            \hline
        \end{tabular}
    \end{center}
    Example:
    \begin{center}
        D'abord, j'\textcolor{red}{ai essayé} de l'appeler au bureau, puis j'\textcolor{red}{ai appelé} son téléphone portable. Enfin, j'\textcolor{red}{ai téléphoné} à sa secrétaire. \\
        (First, I tried calling him at the office, then I called his cell phone. Finally, I called his secretary)
    \end{center}
\end{itemize}

\subsection{Formation}
Passé Composé is formed by 
\begin{center}
    Avoir/Être (conjugated in the present tense) + Past Participle
\end{center}

\subsection{Conjugation of Avoir}

\begin{center}
    \begin{tabular}{|c|c|c|}
        \hline
        \textbf{Pronoun} & \textbf{Conjugation} & \textbf{Pronunciation} \\
        \hline
        je & ai & \textipa{[e]} \\
        tu & as & \textipa{[a]} \\
        il/elle/on & a & \textipa{[a]} \\
        nous & avons & \textipa{[av\~o]} \\
        vous & avez & \textipa{[ave]} \\
        ils/elles & ont & \textipa{[o]} \\
        \hline
    \end{tabular}
\end{center}

\subsection{Conjugation of Être}

\begin{center}
    \begin{tabular}{|c|c|c|}
        \hline
        \textbf{Pronoun} & \textbf{Conjugation} & \textbf{Pronunciation} \\
        \hline
        je & suis & \textipa{[s\~i]} \\
        tu & es & \textipa{[e]} \\
        il/elle/on & est & \textipa{[e]} \\
        nous & sommes & \textipa{[s\~om]} \\
        vous & êtes & \textipa{[et]} \\
        ils/elles & sont & \textipa{[s\~o]} \\
        \hline
    \end{tabular}
\end{center}

\subsection{Participes Passés Réguliers}

\begin{enumerate}
    \item -ER verbs: remove the -ER and add -é, e.g. \textit{manger} $\rightarrow$ \textit{mangé}
    \item -IR verbs: remove the -IR and add -i, e.g. \textit{finir} $\rightarrow$ \textit{fini}
    \item -RE verbs: remove the -RE and add -u, e.g. \textit{vendre} $\rightarrow$ \textit{vendu}
\end{enumerate}

\subsection{Participes Passés Irréguliers}
\begin{itemize}
    \item Avoir $\rightarrow$ eu (to have)
    \item Conduire $\rightarrow$ conduit (to drive)
    \item Être $\rightarrow$ été (to be)
    \item Ouvrir $\rightarrow$ ouvert (to open)
    \item Pleuvoir $\rightarrow$ plu (to rain)
    \item Venir $\rightarrow$ venu (to come)
\end{itemize}

\subsection{Remarques}

\begin{enumerate}
    \item Study by group, for instance:
    \begin{itemize}
        \item Venir $\rightarrow$ v\textcolor{red}{enu}
        \item Tenir $\rightarrow$ t\textcolor{red}{enu} (to hold)
        \item Appart\textcolor{red}{enir} $\rightarrow$ appart\textcolor{red}{enu} (to belong)
        \item Surv\textcolor{red}{enir} $\rightarrow$ surv\textcolor{red}{enu} (to survive)
    \end{itemize}
    and
    \begin{itemize}
        \item Ouvrir $\rightarrow$ ouv\textcolor{red}{ert}
        \item Offrir $\rightarrow$ off\textcolor{red}{ert} (to offer)
    \end{itemize}
    \item Memorize spelling, it will lead to pronunciation change, for example when changing to feminine, we have:
    \begin{itemize}
        \item Conduire: Conduit \textipa{[k\~o.d\~i]} $\rightarrow$ Conduite \textipa{[k\~o.d\~it]}
        \item Ouvrir: Ouvert \textipa{[u.v\~e]} $\rightarrow$ Ouverte \textipa{[u.v\~\textepsilon]}
        \item Prendre: Pris \textipa{[p\~Ri]} $\rightarrow$ Prise \textipa{[p\~Riz]}
        \item Venir: Venu \textipa{[v\~\textepsilon.ny]} $\rightarrow$ Venue \textipa{[v\~\textepsilon.ny]}
    \end{itemize}
    \item The word \textit{déjà} (already) is always between the auxiliary verb and the past participle, for example:
    \begin{center}
        J'ai déjà mangé (I have already eaten)
    \end{center}
\end{enumerate}

\subsection{Passé Composé with Avoir}
In regular conjugation with \textit{avoir}, the past participle does not \textbf{agree in gender and number} with the subject. It agrees with the \textbf{direct object} if it is placed before the verb. 
\begin{center}
    Direct object / Direct object pronoun + avoir + past participle + e.s.es
\end{center}
We can find the direct object by asking Qui? or Quoi? For example:
\begin{center}
    Beefy mange \textcolor{red}{une pomme} $\to$ Beefy \textcolor{red}{la} mange
\end{center}
\begin{itemize}
    \item Beefy mange quoi? Une pomme, so "pomme" is the direct object. The direct object or direct object pronoun is before the verb, so the past participle agrees with it.
\end{itemize}
When conjugation with Avoir agrees with gender and number:
\begin{enumerate}
    \item Direct Object / Direct Object Pronoun
    
    \begin{enumerate}
        \item Feminine COD with no pronounciation change:
        \begin{center}
            Elle a envoyé une letter (She sent a letter) \\
            Elle a envoyé quoi? (What did she send?) \\
            Une lettre 
        \end{center}
        The COD is \textit{une lettre} and DOP is \textit{La / L'}, so the sentence becomes
        \begin{center}
            Elle \textcolor{red}{l'}a envoyé\textcolor{red}{e} (She sent it)
        \end{center}

        \item Feminine COD with pronounciation change:
        \begin{center}
            J'ai écrit une lettre (I wrote a letter) \\
            J'ai écrit quoi? (What did I write?) \\
            Une lettre
        \end{center}
        The COD is \textit{une lettre} and DOP is \textit{La / L'}, so the sentence becomes
        \begin{center}
            Je \textcolor{red}{l'}ai écrit\textcolor{red}{e} (I wrote it)
        \end{center}

        \item Plural COD:
        \begin{center}
            Nous avons acheté des fleurs (We bought flowers) \\
            Nous avons acheté quoi? (What did we buy?) \\
            Des fleurs
        \end{center}
        The COD is \textit{des fleurs} and DOP is \textit{Les}, so the sentence becomes
        \begin{center}
            Nous \textcolor{red}{les} avons acheté\textcolor{red}{es} (We bought them)
        \end{center}
    \end{enumerate}

    \item Relative Pronoun \textit{Que} / \textit{Qu'}
    \begin{itemize}
        \item The relative pronoun replaces the direct object, just another way to word a sentence. Consider:
        \begin{center}
            Elle a envoyé une lettre $\to$ La lettre \textcolor{red}{qu'}elle a envoyé\textcolor{red}{e} \\
            (She sent a letter $\to$ The letter that she sent)
        \end{center}

        \begin{center}
            J'ai reçu les colis du voisin $\to$ Les colis \textcolor{red}{que} j'ai reçu\textcolor{red}{s} \\
            (I got the neighbor's packages $\to$ The neighbour's packages (that) I got)
        \end{center}
    \end{itemize}
    
    
    
\end{enumerate}

\subsection{Passé Composé with Être}
Most of verbs is conjugated with \textit{avoir}, the verbs conjugated with \textit{être} are:
\begin{center}
    DR \& MRS VANDERTRAMPP
\end{center}
\begin{center}
    \begin{tabular}{|c|c|c|c|}
        \hline
        \textbf{Verb} & \textbf{Past Participle} & \textbf{Pronunciation} & \textbf{English Meaning} \\
        \hline
        Devenir & Devenu & \textipa{[d\~\textepsilon.v\~y]} & to become \\
        \hline
        Revenir & Revenu & \textipa{[R\~\textepsilon.v\~y]} & to come back \\
        \hline
        Monter & Monté & \textipa{[m\~\textopeno.t\~e]} & to go up \\
        \hline
        Rester & Resté & \textipa{[R\~\textepsilon.te]} & to stay \\
        \hline
        Sortir & Sorti & \textipa{[s\~oR.ti]} & to go out \\
        \hline
        Venir & Venu & \textipa{[v\~\textepsilon.ny]} & to come \\
        \hline
        Aller & Allé & \textipa{[a.le]} & to go \\
        \hline
        Naître & Né & \textipa{[ne]} & to be born \\
        \hline
        Descendre & Descendu & \textipa{[d\~\textepsilon.s\~y]} & to go down \\
        \hline
        Entrer & Entré & \textipa{[\~a.t\~e]} & to enter \\
        \hline
        Rentrer & Rentré & \textipa{[R\~\textepsilon.t\~e]} & to come back \\
        \hline
        Tomber & Tombé & \textipa{[t\~\textopeno.be]} & to fall \\
        \hline
        Retourner & Retourné & \textipa{[R\~\textepsilon.tuR.ne]} & to return \\
        \hline
        Arriver & Arrivé & \textipa{[a.ri.ve]} & to arrive \\
        \hline
        Mourir & Mort & \textipa{[m\~oR]} & to die \\
        \hline
        Partir & Parti & \textipa{[paR.ti]} & to leave \\
        \hline
        Passer & Passé & \textipa{[pa.se]} & to pass \\
        \hline
    \end{tabular}
\end{center}
The past participle of verbs conjugated with \textit{être} \textbf{agrees in gender and number} with the subject. For example, consider Devenir (to become):
\begin{itemize}
    \item Je suis devenu(e)
    \item Tu es devenu(e)
    \item Il est devenu
    \item Elle est devenue
    \item On est devenu\textcolor{red}{s} (masc. pl.) / devenu\textcolor{red}{es} (fem. pl.)
    \item Nous sommes devenu\textcolor{red}{s} (masc. pl.) / devenu\textcolor{red}{es} (fem. pl.)
    \item Vous êtes devenu\textcolor{red}{s} (masc. pl.) / devenu\textcolor{red}{es} (fem. pl.)
    \item Ils sont devenu\textcolor{red}{s}
    \item Elles sont devenu\textcolor{red}{es}
\end{itemize}

\subsection{Le Passé Composé pour les Verbes Réfléchis}
The structure is simply:
\begin{center}
    Reflexive pronoun + être + past participle + e.s.es
\end{center}
For example, consider \textit{S'adapter} (to adapt):
\begin{center}
    \begin{tabular}{|c|}
        \hline
        Je \textcolor{red}{me} suis adapté\textcolor{red}{(e)} \\
        \hline
        Tu \textcolor{red}{t'}es adapté\textcolor{red}{(e)} \\
        \hline
        Il \textcolor{red}{s'}est adapté \\
        \hline
        Elle \textcolor{red}{s'}est adaptée \\
        \hline
        On \textcolor{red}{s'}est adapté\textcolor{red}{(e)s} \\
        \hline
        Nous \textcolor{red}{nous} sommes adapté\textcolor{red}{(e)s} \\
        \hline
        Vous \textcolor{red}{vous} êtes adapté\textcolor{red}{(e)s} \\
        \hline
        Ils \textcolor{red}{se} sont adapté\textcolor{red}{s} \\
        \hline
        Elles \textcolor{red}{se} sont adapté\textcolor{red}{es} \\
        \hline
    \end{tabular}
\end{center}
Other verbs:
\begin{itemize}
    \item S'asseoir (to sit down): pronounciation changes
    \begin{center}
        \begin{tabular}{|c|}
            \hline
            Je \textcolor{red}{me} suis assis\textcolor{red}{(e)} \\
            \hline
            Tu \textcolor{red}{t'}es assis\textcolor{red}{(e)} \\
            \hline
            Il \textcolor{red}{s'}est assis \\
            \hline
            Elle \textcolor{red}{s'}est assise \\
            \hline
            On \textcolor{red}{s'}est assis\textcolor{red}{(es)} \\
            \hline
            Nous \textcolor{red}{nous} sommes assis\textcolor{red}{(es)} \\
            \hline
            Vous \textcolor{red}{vous} êtes assis\textcolor{red}{(es)} \\
            \hline
            Ils \textcolor{red}{se} sont assis \\
            \hline
            Elles \textcolor{red}{se} sont assises \\
            \hline
        \end{tabular}
    \end{center}
\end{itemize}

\subsection{*** Reflexive Verb and Direct Object ***}
Consider:
\begin{center}
    Elle \textcolor{red}{s'}est coupé\textcolor{red}{e} (She cut herself) \\
    Elle s'est coupé la main (She cut her hand)
    Elle s'est coupée quoi?
\end{center}
Note the COD in this case is \textit{la main} and DOP is \textit{La / L'}, so the sentence becomes
\begin{center}
    Elle se \textcolor{red}{l'}est coupé\textcolor{red}{e} (She cut it)
\end{center}
The direct object is before the verb, so it will agree in gender and number again, so we have the "e" at the end of \textit{coupé}.

\subsection{La Négation du Passé Composé}
We have, except for \textit{Ne...personne}
\begin{center}
    Pronoun + ne + auxiliary verb + pas + past participle
\end{center}
For \textit{Ne...personne}, we have
\begin{center}
    Pronoun + ne + auxiliary verb + past participle + personne
\end{center}
For example:
\begin{center}
    Je n'ai pas mangé (I didn't eat) \\
    Nous n'avons rien reçu (We didn't receive anything) \\
    Je n'ai vu personne (I didn't see anyone)
\end{center}

\begin{enumerate}
    \item For reflexive pronouns, we have
    \begin{center}
        Pronoun + ne + reflexive pronoun + auxiliary verb + pas + past participle
    \end{center}
    For example:
    \begin{center}
        Je ne me suis pas lavé (I didn't wash myself) \\
        Elle ne s'est rien demandé (She didn't ask herself anything) \\
        Nous ne nous sommes jamais demandé(e)s pourquoi (We never asked ourselves why)
    \end{center}
    
    \item For questions, we have
    \begin{enumerate}
        \item Raising your voice:
        \begin{center}
            Tu \textcolor{red}{es} déjà \textcolor{red}{parti}? (Have you already left?) \\
            Ils \textcolor{red}{ont célebré} son augmentation? (Did they celebrate his promotion?) \\
            Tu \textcolor{red}{t'es lavé} les cheveux? (Did you wash your hair?) 
        \end{center}

        \item Est-ce que:
        \begin{center}
            Est-ce que tu \textcolor{red}{es} déjà \textcolor{red}{parti}? \\
            Est-ce qu'ils \textcolor{red}{ont célebré} son augmentation? \\
            Est-ce que tu \textcolor{red}{t'es lavé} les cheveux? 
        \end{center}

        \item Inversion (used for test, very formal and not used in spoken French):
        \begin{center}
            \textcolor{red}{Es}-tu déjà \textcolor{red}{parti}? \\
            \textcolor{red}{Ont}-ils \textcolor{red}{célebré} son augmentation? \\
            Elle \textcolor{red}{a été} en Anustralie $\to$ \textcolor{red}{A-t}-elle \textcolor{red}{été} en Australie? \\
            \textcolor{red}{T'es}-tu \textcolor{red}{lavé} les cheveux?
        \end{center}
    \end{enumerate}
\end{enumerate}

\newpage

\section{Imparfait}

Always use \textit{Nous} form of the verb in the present tense, remove the -ons, and add the following endings:
\begin{center}
    \begin{tabular}{|c|c|c|}
        \hline
        \textbf{Pronoun} & \textbf{Ending} & \textbf{Pronunciation} \\
        \hline
        je & -ais & \textipa{[e]} \\
        tu & -ais & \textipa{[e]} \\
        il/elle/on & -ait & \textipa{[e]} \\
        nous & -ions & \textipa{[j\~o]} \\
        vous & -iez & \textipa{[je]} \\
        ils/elles & -aient & \textipa{[e]} \\
        \hline
    \end{tabular}
\end{center}

\subsection{Quand Utiliser l'Imparfait}
The imparfait is used to talk about the past, mostly actions and situations \textbf{without specific timeframe}, that happened an \textbf{unspecified amount of times} or \textbf{were in progress when something else happened}.
The following cases can be translated to \textit{Je mangeais} in French:
\begin{enumerate}
    \item The past progressive (I was eating)
    \item The simple past (I ate / I used to eat / I would eat)
\end{enumerate}

\begin{enumerate}
    \item To describe habits and repeated actions in the past without being specific in terms of timing, some common seen words:
    \begin{center}
        \begin{tabular}{|c|c|}
            \hline
            \textbf{French} & \textbf{English} \\
            \hline
            Autrefois & Formerly \\
            \hline
            D'habitude & Usually \\
            \hline
            En général & In general \\
            \hline
            Jamais & Never \\
            \hline
            Le lundi & On Mondays \\
            \hline
        \end{tabular}
    \end{center}
    For example:
    \begin{center}
        \textit{Être}: Autrefois, il \textcolor{red}{était} professor (Formerly, he was a professor) \\
        \textit{Étudier}: Nous n'\textcolor{red}{étudiions} jamais (We never studied)
    \end{center}

    \item To talk about repeated actions in the past that could translate to used to or would:
    \begin{center}
        \textit{Dormir}: Il \textcolor{red}{dormait} souvent avec la fenêtre ouverte (He often used to sleep with the window open)  
    \end{center}

    \item To describe a person, a property, the weather or even a physical state or an emotion:
    \begin{center}
        \textit{Se situner}: La maison \textcolor{red}{se situait} au bout de la rue (The house was situated at the end of the street) \\
        \textit{Pleuvoir}: Il \textcolor{red}{pleuvait} beaucoup (It was raining a lot)
    \end{center}
    Some words include:
    \begin{center}
        \begin{tabular}{|c|c|}
            \hline
            \textbf{French} & \textbf{English} \\
            \hline
            Aimer & To like \\
            \hline
            Croire & To believe \\
            \hline
            Être & To be \\
            \hline
            S'attendre & To expect \\
            \hline
            Souhaiter & To wish \\
            \hline
        \end{tabular}
    \end{center}
    For example:
    \begin{center}
        \textit{Aimer}: J'\textcolor{red}{aimais} beaucoup le café (I used to like coffee a lot) \\
        \textit{Souhaiter}: Il \textcolor{red}{souhaitait} réussir (He wished to succeed)
    \end{center}

    \item To talk about actions happening at the same time:
    \begin{center}
        Elle \textcolor{red}{rangeait} la terrasse pendant que les enfants \textcolor{red}{jouaient} (She was tidying up the patio while the kids were playing) \\
        Il \textcolor{red}{travaillait} et elle \textcolor{red}{répondait} au téléphone (He was working and she was answering the phone)
    \end{center}

    \item Particular cases:
    \begin{enumerate}
        \item After \textit{Si} (if):
        \begin{center}
            Si seulement j'\textcolor{red}{avais} plus de temps! (If only I had more time!)
        \end{center}
        \item After \textit{Et Si} (what if):
        \begin{center}
            Et si tu \textcolor{red}{venais} ce week-end? (What if you came this weekend?) \\
            Et si on \textcolor{red}{se voyait} ce soir? (What about we meet tonight?)
        \end{center}
    \end{enumerate}

    \item Sentence built with infinitive:
    \begin{enumerate}

        \item Aller + infinitive:
        \begin{center}
            J'\textcolor{red}{allais} partir (I was going to leave) \\
            Il \textcolor{red}{allait} la demander en mariage (He was going to ask her to marry him)
        \end{center}

        \item Venir de + infinitive:
        \begin{center}
            Il \textcolor{red}{venait} d'arriver (He had just arrived)
        \end{center}

        \item Être en train de + infinitive (to be in the process of):
        \begin{center}
            Nous \textcolor{red}{étions} en train d'étudier (We were studying)
        \end{center}

        \item Être sur le point de + infinitive (to be about to):
        \begin{center}
            Ils \textcolor{red}{étaient} sur le point de divorcer (They were about to get divorced)
        \end{center}
    \end{enumerate}
\end{enumerate}

\subsection{Irréguliers}
\begin{itemize}
    \item Verbs present \textit{Nous} has -ç, change to -c for \textit{Nous, Vous}, keep for other ones. This is because when 'c' followed by an 'i', it makes the 's' sound already, no need to add the cedilla.
    For example (Prononcer):
    \begin{center}
        Nous prononçons \\
        \begin{tabular}{|c|c|}
            \hline
            Je & pronon\textbf{ç}ais \\
            Tu & pronon\textbf{ç}ais \\
            Il/Elle/On & pronon\textbf{ç}ait \\
            Nous & pronon\textbf{c}ions \\
            Vous & pronon\textbf{c}iez \\
            Ils/Elles & pronon\textbf{ç}aient \\
            \hline
        \end{tabular}
    \end{center}

    \item Étudier (to study): a very special spelling for Imparfait, double 'i' for \textit{Nous, Vous}, keep for other ones, purely for spelling:
    \begin{center}
        Nous étudiions \\
        \begin{tabular}{|c|c|}
            \hline
            Je & étudiais \\
            \hline
            Tu & étudiais \\
            \hline
            Il/Elle/On & étudiait \\
            \hline
            Nous & étud\textbf{i}ions \\
            \hline
            Vous & étud\textbf{i}iez \\
            \hline
            Ils/Elles & étudiaient \\
            \hline
        \end{tabular}
    \end{center}

    \item Être: a very special verb, in Passé Composé, it is \textit{été} and in Imparfait, the stem is \textit{ét}:
    \begin{center}
        \begin{tabular}{|c|c|}
            \hline
            J' & étais \\
            \hline
            Tu & étais \\
            \hline
            Il/Elle/On & était \\
            \hline
            Nous & étions \\
            \hline
            Vous & étiez \\
            \hline
            Ils/Elles & étaient \\
            \hline
        \end{tabular}
    \end{center}
    
    \item Other irregular verbs:
    \begin{center}
        \begin{tabular}{|c|c|c|}
            \hline
            Faire (to do) & Nous \textcolor{red}{fais}ons & Je fais\textcolor{red}{ais} \\
            \hline
            Croire (to believe) & Nous \textcolor{red}{croy}ons & Je croy\textcolor{red}{ais} \\
            \hline
            Connaître (to know) & Nous \textcolor{red}{connaiss}ons & Je connaiss\textcolor{red}{ais} \\
            \hline
            Appeler (to call) & Nous \textcolor{red}{appel}ons & J'appel\textcolor{red}{ais} \\
            \hline
            Répéter (to repeat) & Nous \textcolor{red}{répét}ons & Je répét\textcolor{red}{ais} \\
            \hline
            Prendre (to take) & Nous \textcolor{red}{pren}ons & Je pren\textcolor{red}{ais} \\
            \hline
            Conduire (to drive) & Nous \textcolor{red}{conduis}ons & Je conduis\textcolor{red}{ais} \\
            \hline
            Dire (to say) & Nous \textcolor{red}{dis}ons & Je dis\textcolor{red}{ais} \\
            \hline
        \end{tabular}
    \end{center}
\end{itemize}

\newpage

\section{Quand Imparfait et Passé Composé dans la Même Phrase}
Consider the following:
\begin{enumerate}
    \item 
    \begin{center}
        Je \textcolor{red}{regardais} la télévision quand le téléphone \textcolor{red}{a sonné} (I was watching TV when the phone rang)
    \end{center}
    The Imparfait is \textit{regardais} and the Passé Composé is \textit{a sonné}. The phone rang at the exact moment when I was watching TV. We can pin point the exact moment when the phone rang.

    \item 
    \begin{center}
        Elle \textcolor{red}{mangeait} quand sa soeur \textcolor{red}{est arrivée} (She was eating when her sister arrived)
    \end{center}
    The Imparfait is \textit{mangeait} and the Passé Composé is \textit{est arrivée}. The sister arrived at the exact moment when she was eating. We can pin point the exact moment when the sister arrived.
    
\end{enumerate}

\section{Le Passé Récent}
Something happened recently. We conjugate \textit{Venir} in present tense 
\begin{center}
    \begin{tabular}{|c|}
        \hline
        Je \textcolor{red}{viens} \\
        \hline
        Tu \textcolor{red}{viens} \\
        \hline
        Il/Elle/On \textcolor{red}{vient} \\
        \hline
        Nous \textcolor{red}{venons} \\
        \hline
        Vous \textcolor{red}{venez} \\
        \hline
        Ils/Elles \textcolor{red}{viennent} \\
        \hline
    \end{tabular}
\end{center}
then
\begin{center}
    Venir + de + infinitive verb
\end{center}
For example:
\begin{center}
    Je \textcolor{red}{viens de} + infinitive (I have just done / I just did)
\end{center}
We can also add \textit{juste} in between, it's in front of the verb in French,
\begin{center}
    Je \textcolor{red}{viens juste de} + infinitive (I have just done / I just did)
\end{center}
Do not mix with 
\begin{center}
    Venir de + location (to be from)
\end{center}
Example sentences:
\begin{center}
    Marie \textcolor{red}{vient de} m'appeler (Marie just called me) \\
    Nous \textcolor{red}{venons d'}arriver (We just arrived)
\end{center}

\subsection{Le Passé Récent in the past}
Something happened just before something else happened in the past. We conjugate \textit{Venir} in Imparfait tense, then
\begin{center}
    Je \textcolor{red}{venais de} + infinitive (I had just done) \\
    Je \textcolor{red}{venais juste de} + infinitive (I had just done)
\end{center}
Venir in Imparfait:
\begin{center}
    \begin{tabular}{|c|}
        \hline
        Je \textcolor{red}{venais} \\
        \hline
        Tu \textcolor{red}{venais} \\
        \hline
        Il/Elle/On \textcolor{red}{venait} \\
        \hline
        Nous \textcolor{red}{venions} \\
        \hline
        Vous \textcolor{red}{veniez} \\
        \hline
        Ils/Elles \textcolor{red}{vienaient} \\
        \hline
    \end{tabular}
\end{center}
Note the second verb is in Passé Composé, for example:
\begin{center}
    Je \textcolor{red}{venais juste de} traverser la rue quand l'ambulance \textcolor{red}{est passée} \\
    (I had just crossed the street when the ambulance passed by) \\
    Elle \textcolor{red}{venait juste d'}acheter cette voiture quand elle \textcolor{red}{est tombée en panne} \\
    (She had just bought this car when it broke down)
\end{center}

\subsection{Negation with Le Passé Récent}
Does not exist.

\chapter{Future}

\section{Le Futur Simple}
This is \textit{will} in English, used to talk about future plans and intentions. The stem used for the \textbf{futur simple} will be used for \textbf{conditionnel} as well. \\
Some keywords are:
\begin{center}
    \begin{tabular}{|c|c|}
        \hline
        \textbf{French} & \textbf{English} \\
        \hline
        Demain & Tomorrow \\
        \hline
        Plus tard & Later \\
        \hline
        Bientôt & Soon \\
        \hline
        Un jour & One day \\
        \hline
        Ce soir & Tonight \\
        \hline
        La semaine prochaine & Next week \\
        \hline
        L'année prochaine & Next year \\
        \hline
    \end{tabular}
\end{center}

\begin{enumerate}
    \item To simply talk about something that is going to happen in the future:
    \begin{center}
        Un jour, je \textcolor{red}{visiterai} l'Australie (One day, I will visit Australia)
    \end{center}

    \item To make suppositions about the future:
    \begin{center}
        Je pense qu'il \textcolor{red}{arrivera} en retard (I think he will arrive late)
    \end{center}

    \item ***After conjunctions of time (different from english), some reference are:
    \begin{center}
        \begin{tabular}{|c|c|}
            \hline
            \textbf{French} & \textbf{English} \\
            \hline
            Quand & When \\
            \hline
            Lorsque / Lorsqu' & When \\
            \hline
            Dès que / Dès qu' & As soon as \\
            \hline
            Aprés que / Aprés qu' & After \\
            \hline
        \end{tabular}
    \end{center}
    For example;
    \begin{center}
        Appelle-moi quand tu \textcolor{red}{recevras} l'invitation \\
        (Call me when you get the invitation) \\
        Vois avec elle dès qu'elle \textcolor{red}{sera} de retour \\
        (See that with her when she is back) 
    \end{center}

\end{enumerate}

\subsection{-ER Verbs}
\begin{enumerate}
    \item The infinity verb follow by the stems, consider \textit{Ranger} (to tidy):
        \begin{center}
            \begin{tabular}{|c|c|}
                \hline
                \textbf{Conjugation} & \textbf{Pronunciation} \\
                \hline
                Je ranger\textbf{ai} & \textipa{[e]} \\
                \hline
                Tu ranger\textbf{as} & \textipa{[a]} \\
                \hline
                Il/Elle/On ranger\textbf{a} & \textipa{[a]} \\
                \hline
                Nous ranger\textbf{ons} & \textipa{[j\~o]} \\
                \hline
                Vous ranger\textbf{ez} & \textipa{[e]} \\
                \hline
                Ils/Elles ranger\textbf{ont} & \textipa{[o]} \\
                \hline
            \end{tabular}
        \end{center}
    \item The verbs that keeps the same stem as they are in infinitives are:
        \begin{enumerate}
            \item -cer: Lan\textcolor{red}{cer} $\rightarrow$ Je lancer\textcolor{red}{ai} (to throw)
            \item -ger: Man\textcolor{red}{ger} $\rightarrow$ Je manger\textcolor{red}{ai} (to eat)
            \item -éder: C\textcolor{red}{éder} $\rightarrow$ Je céder\textcolor{red}{ai} (to give up)
            \item -érer: Esp\textcolor{red}{érer} $\rightarrow$ J'espérer\textcolor{red}{ai} (to hope)
            \item -éter: Rép\textcolor{red}{éter} $\rightarrow$ Je répéter\textcolor{red}{ai} (to repeat)
            \item -ébrer: Cél\textcolor{red}{ébrer} $\rightarrow$ Je célébrer\textcolor{red}{ai} (to celebrate)
        \end{enumerate}
    \item The verbs change the stem are:
        \begin{enumerate}
            \item -yer: 
            \begin{itemize}
                \item Netto\textcolor{red}{yer} $\rightarrow$ Je netto\textcolor{red}{i}er\textcolor{red}{ai} (to clean)
                \item Essa\textcolor{red}{yer} $\rightarrow$ J'essa\textcolor{red}{i}er\textcolor{red}{ai} (to try)
                \item Essu\textcolor{red}{yer} $\rightarrow$ J'essu\textcolor{red}{i}er\textcolor{red}{ai} (to wipe)
            \end{itemize}
            \item Appeler (to call), same for Jeter $\to$ Je jetter\textcolor{red}{ai}
            \begin{center}
                \begin{tabular}{|c|c|}
                    \hline
                    J'appeller\textcolor{red}{ai} & \textipa{[a.p\~\textepsilon.le]} \\
                    \hline
                    Tu appeller\textcolor{red}{as} & \textipa{[a.p\~\textepsilon.le]} \\
                    \hline
                    Il/Elle/On appeller\textcolor{red}{a} & \textipa{[a.p\~\textepsilon.le]} \\
                    \hline
                    Nous appeller\textcolor{red}{ons} & \textipa{[a.p\~\textepsilon.l\~o]} \\
                    \hline
                    Vous appeller\textcolor{red}{ez} & \textipa{[a.p\~\textepsilon.le]} \\
                    \hline
                    Ils/Elles appeller\textcolor{red}{ont} & \textipa{[a.p\~\textepsilon.le]} \\
                    \hline
                \end{tabular}
            \end{center}
        \end{enumerate}
    \item e changes to è in the stem:
    \begin{itemize}
        \item G\textcolor{red}{eler} $\rightarrow$ Je g\textcolor{red}{è}ler\textcolor{red}{ai} (to freeze)
        \item Am\textcolor{red}{ener} $\rightarrow$ J'am\textcolor{red}{è}ner\textcolor{red}{ai} (to bring)
        \item P\textcolor{red}{eser} $\rightarrow$ Je p\textcolor{red}{è}ser\textcolor{red}{ai} (to weigh)
        \item Ach\textcolor{red}{eter} $\rightarrow$ J'ach\textcolor{red}{è}ter\textcolor{red}{ai} (to buy)
        \item Ach\textcolor{red}{ever} $\rightarrow$ J'ach\textcolor{red}{è}ver\textcolor{red}{ai} (to finish)
    \end{itemize}
\end{enumerate}

\subsection{-IR Verbs}
Same as -ER verbs, the stem is the same as the infinitive, consider \textit{Accomplir} (to accomplish):
\begin{center}
    \begin{tabular}{|c|c|}
        \hline
        \textbf{Conjugation} & \textbf{Pronunciation} \\
        \hline
        J'accomplir\textbf{ai} & \textipa{[e]} \\
        \hline
        Tu accomplir\textbf{as} & \textipa{[a]} \\
        \hline
        Il/Elle/On accomplir\textbf{a} & \textipa{[a]} \\
        \hline
        Nous accomplir\textbf{ons} & \textipa{[j\~o]} \\
        \hline
        Vous accomplir\textbf{ez} & \textipa{[e]} \\
        \hline
        Ils/Elles accomplir\textbf{ont} & \textipa{[o]} \\
        \hline
    \end{tabular}
\end{center}
The word \textit{Accueill\textcolor{red}{ir}} is
\begin{center}
    J'accueill\textcolor{red}{er}ai, \textipa{[a.k\~\textepsilon.je.le]}
\end{center}

\subsection{-RE Verbs}
Remove the -e, consider \textit{Prendre} (to take):
\begin{center}
    \begin{tabular}{|c|c|}
        \hline
        \textbf{Conjugation} & \textbf{Pronunciation} \\
        \hline
        Je prendr\textbf{ai} & \textipa{[e]} \\
        \hline
        Tu prendr\textbf{as} & \textipa{[a]} \\
        \hline
        Il/Elle/On prendr\textbf{a} & \textipa{[a]} \\
        \hline
        Nous prendr\textbf{ons} & \textipa{[j\~o]} \\
        \hline
        Vous prendr\textbf{ez} & \textipa{[e]} \\
        \hline
        Ils/Elles prendr\textbf{ont} & \textipa{[o]} \\
        \hline
    \end{tabular}
\end{center}
Other verbs include:
\begin{itemize}
    \item Prendre (to take): Je prendrai
    \item Craindre (to fear): Je craindrai
    \item Éteindre (to turn off): J'éteindrai
    \item Joindre (to join): Je joindrai
    \item Battre (to beat): Je battrai
\end{itemize}

\subsection{Irregular Verbs}
These irregular verbs are the same for \textbf{conditionnel} as well.
\begin{center}
    \begin{tabular}{|c|c|}
        \hline
        \textbf{Verb} & \textbf{Future Simple} \\
        \hline
        Avoir (to have) & J'aur\textcolor{red}{ai} \\
        \hline
        Être (to be) & Je ser\textcolor{red}{ai} \\
        \hline
        Faire (to do) & Je fer\textcolor{red}{ai} \\
        \hline
        Pouvoir (to be able to) & Je pourr\textcolor{red}{ai} \\
        \hline
        Savoir (to know) & Je saur\textcolor{red}{ai} \\
        \hline
        Venir (to come) & Je viendr\textcolor{red}{ai} \\
        \hline
        Voir (to see) & Je verr\textcolor{red}{ai} \\
        \hline
        Vouloir (to want) & Je voudr\textcolor{red}{ai} \\
        \hline
        Aller (to go) & J'ir\textcolor{red}{ai} \\
        \hline
        Devoir (to have to) & Je devr\textcolor{red}{ai} \\
        \hline
        Recevoir (to receive) & Je recevr\textcolor{red}{ai} \\
        \hline
        Envoyer (to send) & J'enverr\textcolor{red}{ai} \\
        \hline
        Mourir (to die) & Je mourr\textcolor{red}{ai} \\
        \hline
        Courir (to run) & Je courr\textcolor{red}{ai} \\
        \hline
        Plaire (to please) & Je plair\textcolor{red}{ai} \\
        \hline
        Pleuvoir (to rain) & Il pleuvr\textcolor{red}{a} \\
        \hline
        Falloir (to be necessary) & Il faudr\textcolor{red}{a} \\
        \hline
        Valoir (to be worth) & Il vaudr\textcolor{red}{a} \\
        \hline
    \end{tabular}
\end{center}

\section{Le Futur Proche}
Use \textit{Aller} in present tense:
\begin{center}
    Aller + infinitive verb
\end{center}
for example:
\begin{center}
    Je \textbf{vais} + infinitive (I am going to)
\end{center}

\subsection{La Conjugation de Aller}
\begin{center}
    \begin{tabular}{|c|c|}
        \hline
        \textbf{Conjugation} & \textbf{Pronunciation} \\
        \hline
        Je \textcolor{red}{vais} & \textipa{[v\textepsilon]} \\
        \hline
        Tu \textcolor{red}{vas} & \textipa{[va]} \\
        \hline
        Il/Elle/On \textcolor{red}{va} & \textipa{[va]} \\
        \hline
        Nous \textcolor{red}{allons} & \textipa{[a.l\~o]} \\
        \hline
        Vous \textcolor{red}{allez} & \textipa{[a.le]} \\
        \hline
        Ils/Elles \textcolor{red}{vont} & \textipa{[v\~o]} \\
        \hline
    \end{tabular}
\end{center}
Example sentences:
\begin{center}
    Elles \textbf{vont apprendre} la cuisine (They are going to learn how to cook) 
\end{center}

\subsection{Reflexive Verbs}
When we have reflexive verb, it is going to agree to the subject, consider:
\begin{center}
    Je \textbf{vais \underline{me} faire} un café (I am going to make myself a coffee) \\
    Tu \textbf{vas \underline{te} faire} un café (You are going to make yourself a coffee)
\end{center}

\subsection{Négation avec Le Futur Proche}

\begin{enumerate}
    \item Simple negation:
        \begin{center}
            Je + ne + vais + pas + infinitive
        \end{center}
        For example:
        \begin{center}
            On ne va pas nager (We are not going to swim)
        \end{center}

    \item Negation with reflexive verb:
        \begin{center}
            Je + ne + vais + pas + infinitive RV
        \end{center}
        For example:
        \begin{center}
            Je vais me faire un café (I am going to make myself a coffee) \\
            Je ne vais pas me faire un café (I am not going to make myself a coffee)
        \end{center}
\end{enumerate}

\chapter{Conditionnel}

\section{Le Conditionnel Présent}
In english it is \textit{would}, use to ask something politely, to express a wish or to talk about something that would happen \textit{if}, about a possibility. \\
To form the conditionnel présent, \textbf{add the endings of the imparfait to the stem used with the futur simple}. \\
Example sentences:
\begin{center}
    Possibility: On \textbf{pourrait} y aller demain (We could go there tomorrow) \\
    Polite: \textbf{Pourrais}-tu me donner son numéro? (Could you give me his number?) \\
    Polite: \textbf{Sauriez}-vous où je peux trouver un bon restaurant? (Would you know where I can find a good restaurant?) \\
    Wish: J'\textbf{souhaiterais} gagner à la loterie (I would like to win the lottery) \\
    Wish: \textbf{Aimerais}-tu aller au cinéma? (Would you like to go to the theater?)
\end{center}
Common verbs used with wish conditionnel:
\begin{center}
    \begin{tabular}{|c|c|}
        \hline
        \textbf{French} & \textbf{English} \\
        \hline
        Aimer & To like \\
        \hline
        Vouloir & To want \\
        \hline
        Désirer & To desire \\
        \hline
        Souhaiter & To wish \\
        \hline
    \end{tabular}
\end{center}

\subsection{Hypothetical Situation}
Usually used with \textit{Si} (if), we have two possible structures:
\begin{enumerate}
    \item Condition first:
        \begin{center}
            Si + Imparfait, Conditionnel Présent
        \end{center}
        \textit{Si} at the beginning, is always going to be followed by Imparfait, the second part of the sentence is going to be Conditionnel Présent. For example:
        \begin{center}
            Si je \textbf{gagnais} à la loterie, j'\textbf{achèterais} une maison (If I won the lottery, I would buy a house) 
        \end{center}
    \item Condition second:
        \begin{center}
            Conditionnel Présent, Si + Imparfait
        \end{center}
        The above example is equivalent to:
        \begin{center}
            J'\textbf{achèterais} une maison si je \textbf{gagnais} à la loterie (I would buy a house, if I won the lottery)
        \end{center}
\end{enumerate}


\subsection{-ER and -IR Verbs}
The endings are the same as imparfait, consider \textit{Ranger} (to tidy):
\begin{center}
    \begin{tabular}{|c|c|}
        \hline
        \textbf{Conjugation} & \textbf{Pronunciation} \\
        \hline
        Je ranger\textcolor{red}{ais} & \textipa{[e]} \\
        \hline
        Tu ranger\textcolor{red}{ais} & \textipa{[e]} \\
        \hline
        Il/Elle/On ranger\textcolor{red}{ait} & \textipa{[e]} \\
        \hline
        Nous ranger\textcolor{red}{ions} & \textipa{[j\~o]} \\
        \hline
        Vous ranger\textcolor{red}{iez} & \textipa{[je]} \\
        \hline
        Ils/Elles ranger\textcolor{red}{aient} & \textipa{[e]} \\
        \hline
    \end{tabular}
\end{center}
Words that doesn't have spelling change, same as imparfait:
\begin{center}
    \begin{tabular}{|c|c|c|}
        \hline
        \textbf{Verb} & \textbf{Conditionnel Présent} & \textbf{English Meaning} \\
        \hline
        Lan\textcolor{red}{cer} & Je lancer\textcolor{red}{ais} & I would throw \\
        \hline
        Man\textcolor{red}{ger} & Je manger\textcolor{red}{ais} & I would eat \\
        \hline
        C\textcolor{red}{éder} & Je céder\textcolor{red}{ais} & I would give up \\
        \hline
        Esp\textcolor{red}{érer} & J'espérer\textcolor{red}{ais} & I would hope \\
        \hline
        Rép\textcolor{red}{éter} & Je répéter\textcolor{red}{ais} & I would repeat \\
        \hline
        Cél\textcolor{red}{ébrer} & Je célébrer\textcolor{red}{ais} & I would celebrate \\
        \hline
    \end{tabular}
\end{center}
Verb changes spelling, same as imparfait:
\begin{center}
    \begin{tabular}{|c|c|c|}
        \hline
        \textbf{Verb} & \textbf{Conditionnel Présent} & \textbf{English Meaning} \\
        \hline
        Netto\textcolor{red}{yer} & Je netto\textcolor{red}{i}er\textcolor{red}{ais} & I would clean \\
        \hline
        Essa\textcolor{red}{yer} & J'essa\textcolor{red}{i}er\textcolor{red}{ais} & I would try \\
        \hline
        Essu\textcolor{red}{yer} & J'essu\textcolor{red}{i}er\textcolor{red}{ais} & I would wipe \\
        \hline
    \end{tabular}
\end{center}
\textit{Appeler} (to call), double 'l', same for \textit{Jeter}:
\begin{center}
    \begin{tabular}{|c|}
        \hline
        J'appeller\textcolor{red}{ais} \\
        \hline
        Tu appeller\textcolor{red}{ais} \\
        \hline
        Il/Elle/On appeller\textcolor{red}{ait} \\
        \hline
        Nous appeller\textcolor{red}{ions} \\
        \hline
        Vous appeller\textcolor{red}{iez} \\
        \hline
        Ils/Elles appeller\textcolor{red}{aient} \\
        \hline
    \end{tabular}
\end{center}
Words don't have accent, same as imparfait:
\begin{center}
    \begin{tabular}{|c|c|c|}
        \hline
        \textbf{Verb} & \textbf{Conditionnel Présent} & \textbf{English Meaning} \\
        \hline
        G\textcolor{red}{eler} & Je g\textcolor{red}{è}ler\textcolor{red}{ais} & I would freeze \\
        \hline
        Am\textcolor{red}{ener} & J'am\textcolor{red}{è}ner\textcolor{red}{ais} & I would bring \\
        \hline
        P\textcolor{red}{eser} & Je p\textcolor{red}{è}ser\textcolor{red}{ais} & I would weigh \\
        \hline
        Ach\textcolor{red}{eter} & J'ach\textcolor{red}{è}ter\textcolor{red}{ais} & I would buy \\
        \hline
        Ach\textcolor{red}{ever} & J'ach\textcolor{red}{è}ver\textcolor{red}{ais} & I would finish \\
        \hline
    \end{tabular}
\end{center}

\subsection{-RE Verbs}
Remove the -e, consider \textit{Prendre} (to take):
\begin{center}
    \begin{tabular}{|c|c|}
        \hline
        \textbf{Conjugation} & \textbf{Pronunciation} \\
        \hline
        Je prendr\textcolor{red}{ais} & \textipa{[e]} \\
        \hline
        Tu prendr\textcolor{red}{ais} & \textipa{[e]} \\
        \hline
        Il/Elle/On prendr\textcolor{red}{ait} & \textipa{[e]} \\
        \hline
        Nous prendr\textcolor{red}{ions} & \textipa{[j\~o]} \\
        \hline
        Vous prendr\textcolor{red}{iez} & \textipa{[je]} \\
        \hline
        Ils/Elles prendr\textcolor{red}{aient} & \textipa{[e]} \\
        \hline
    \end{tabular}
\end{center}

\subsection{Irregular Verbs}
Same as futur simple, we have:
\begin{center}
    \begin{tabular}{|c|c|}
        \hline
        \textbf{Verb} & \textbf{Conditionnel Présent} \\
        \hline
        Avoir (to have) & J'aur\textcolor{red}{ais} \\
        \hline
        Être (to be) & Je ser\textcolor{red}{ais} \\
        \hline
        Faire (to do) & Je fer\textcolor{red}{ais} \\
        \hline
        Pouvoir (to be able to) & Je pourr\textcolor{red}{ais} \\
        \hline
        Savoir (to know) & Je saur\textcolor{red}{ais} \\
        \hline
        Venir (to come) & Je viendr\textcolor{red}{ais} \\
        \hline
        Voir (to see) & Je verr\textcolor{red}{ais} \\
        \hline
        Vouloir (to want) & Je voudr\textcolor{red}{ais} \\
        \hline
        Aller (to go) & J'ir\textcolor{red}{ais} \\
        \hline
        Devoir (to have to) & Je devr\textcolor{red}{ais} \\
        \hline
        Recevoir (to receive) & Je recevr\textcolor{red}{ais} \\
        \hline
        Envoyer (to send) & J'enverr\textcolor{red}{ais} \\
        \hline
        Mourir (to die) & Je mourr\textcolor{red}{ais} \\
        \hline
        Courir (to run) & Je courr\textcolor{red}{ais} \\
        \hline
        Plaire (to please) & Je plair\textcolor{red}{ais} \\
        \hline
        Pleuvoir (to rain) & Il pleuvr\textcolor{red}{ait} \\
        \hline
        Falloir (to be necessary) & Il faudr\textcolor{red}{ait} \\
        \hline
        Valoir (to be worth) & Il vaudr\textcolor{red}{ait} \\
        \hline
    \end{tabular}
\end{center}

\subsection{Devoir - Vouloir - Pouvoir}
These verbs are usually followed by an infinitive verb:
\begin{center}
    \begin{tabular}{|c|c|c|}
        \hline
        \textbf{Verb} & \textbf{Conditionnel Présent} & \textbf{English Meaning} \\
        \hline
        Devoir (should) & Je devr\textcolor{red}{ais} & I would have to \\
        \hline
        Vouloir (would like) & Je voudr\textcolor{red}{ais} & I would want \\
        \hline
        Pouvoir (could) & Je pourr\textcolor{red}{ais} & I would be able to \\
        \hline
    \end{tabular}
\end{center}
Note the conditionnel \textit{pouvoir} is used to ask things very politely, for example:
\begin{center}
    Nous \textbf{pourrions} peut-être aider (We could maybe help) \\
    \textbf{Pourriez}-vous vous taire? (Could you be quiet?)
\end{center}

\chapter{L'Impératif}
Use to give order and advise, only used with \textbf{tu, nous, vous}. Ignore the subject pronoun. 
The three groups of verbs are:
\begin{enumerate}
    \item Regarder (to look):
    \begin{center}
        \begin{tabular}{|c|c|}
            \hline
            \textbf{Present Tense} & \textbf{Impératif} \\
            \hline
            Tu regarde\textcolor{red}{s} & Regarde \\
            \hline
            Nous regardons & Regardons \\
            \hline
            Vous regardez & Regardez \\
            \hline
        \end{tabular}
    \end{center}
    Note for ending -es only, we remove the -s, except for when it is followed by \textit{en, y}, for example:
    \begin{center}
        Commandes-en (Order some) \\
        Restes-y (Stay there)
    \end{center}
    
    \item Finir (to finish):
    \begin{center}
        \begin{tabular}{|c|c|}
            \hline
            \textbf{Present Tense} & \textbf{Impératif} \\
            \hline
            Tu finis & Finis \\
            \hline
            Nous finissons & Finissons \\
            \hline
            Vous finissez & Finissez \\
            \hline
        \end{tabular}
    \end{center}

    \item Attendre (to wait):
    \begin{center}
        \begin{tabular}{|c|c|}
            \hline
            \textbf{Present Tense} & \textbf{Impératif} \\
            \hline
            Tu attends & Attends \\
            \hline
            Nous attendons & Attendons \\
            \hline
            Vous attendez & Attendez \\
            \hline
        \end{tabular}
    \end{center}
\end{enumerate}

\section{Irregular Verbs}
The four irregular verbs are:
\begin{enumerate}
    \item Avoir (to have):
    \begin{center}
        \begin{tabular}{|c|}
            \hline
            Aie \\
            \hline
            Ayons \\
            \hline
            Ayez \\
            \hline
        \end{tabular}
    \end{center}

    \item Être (to be):
    \begin{center}
        \begin{tabular}{|c|}
            \hline
            Sois \\
            \hline
            Soyons \\
            \hline
            Soyez \\
            \hline
        \end{tabular}
    \end{center}

    \item Savoir (to know):
    \begin{center}
        \begin{tabular}{|c|}
            \hline
            Sache \\
            \hline
            Sachons \\
            \hline
            Sachez \\
            \hline
        \end{tabular}
    \end{center}

    \item Vouloir (to want), we only use \textit{veuillez}, the first two exist for nothing:
    \begin{center}
        \begin{tabular}{|c|}
            \hline
            \textit{Veuille} \\
            \hline
            \textit{Veuilons} \\
            \hline
            \textbf{Voulez} \\
            \hline
        \end{tabular}
    \end{center}
\end{enumerate}

\section{Negation with L'Impératif}
The format is:
\begin{center}
    Ne + verb + pas
\end{center}
for example:
\begin{center}
    Ne regarde pas (Don't look) \\
    Ne partez pas (Don't leave) \\
    Ne bouge plus (Don't move anymore / stop moving)
\end{center}

\subsection{Remarques}
For verbs that are followed by preposition (i.e., à, de etc.), we need to include the preposition in the negation, for example:
\begin{center}
    Ne touche \textcolor{red}{à} rien (Don't touch anything)
\end{center}

\section{Object Pronouns with L'Impératif}
Always after the verb, joined by a hyphen. Some examples are:
\begin{center}
    \begin{tabular}{|c|c|}
        \hline
        Demande-moi & Ask me \\
        Regarde-toi & Look at you \\
        Demandez-lui & Ask him / her \\
        Regardez-les & Look at them \\
        Demande-nous & Ask us \\
        Allons-y & Let's go (y refers to place) \\
        Prends-en & Take some (en refers to quantity) \\
        \hline
    \end{tabular}
\end{center}

\subsection{Negation}
Moi changes back to me, toi changes back to te, for example:
\begin{center}
    Ne me demande pas (Don't ask me) \\
    Ne te regarde pas (Don't look at you)
\end{center}

\subsection{DOP and IOP Together}
\begin{enumerate}
    \item In informative sentences, after the verb, we will have:
        \begin{center}
            \begin{tabular}{|c|c|c|c|}
                \hline
                Le / L' & Moi / M' & Nous & Y \\
                La / L' & Toi / T' & Vous & En \\
                Les & & Lui & \\
                & & Leur & \\
                \hline
            \end{tabular}
        \end{center}
        For example:
        \begin{center}
            Donne le jouet au chien (Give the toy to the dog) \\
            Donne-le au chien (Give it to the dog) \\
            Donne-le-lui (Give it to him)
        \end{center}

    \item Negation, the order changes, before the verb:
        \begin{center}
            \begin{tabular}{|c|c|c|c|}
                \hline
                Me / M' & Le / L' & Lui & Y \\
                Te / T' & La / L' & Leur & En \\
                Nous & Les & & \\
                Vous & & & \\
                \hline
            \end{tabular}
        \end{center}
        For example:
        \begin{center}
            Ne le donne pas au chien (Don't give it to the dog) \\
            Ne le lui donne pas (Don't give it to him) \\
            Ne t'en fais pas (Don't worry) \\
            Ne m'y rejoignez pas (Don't meet me there) \\
            Ne l'appelle pas (Don't call him / her)
        \end{center}
\end{enumerate}

\subsection{Direct Object Pronouns}
Replaces the noun that receives the action of the verb directly.
\begin{center}
    \begin{tabular}{|c|c|}
        \hline
        \textbf{Singular} & \textbf{Plural} \\
        \hline
        Me / M' & Nous \\
        \hline
        Te / T' & Vous \\
        \hline
        Le / L' & Les \\
        \hline
        La / L' & \\
        \hline
    \end{tabular}
\end{center}
For \textit{Me, Te}, they are going to become \textit{Moi, Toi} when they are placed after the verb (in informative sentences). Can be seen as we don't usually have a [e] sound at the end of the verb.

\subsection{Indirect Object Pronouns}
Replaces the noun that receives the action of the verb indirectly, usually followed by \textit{à}.
\begin{center}
    \begin{tabular}{|c|c|}
        \hline
        \textbf{Singular} & \textbf{Plural} \\
        \hline
        Me / M' & Nous \\
        \hline
        Te / T' & Vous \\
        \hline
        Lui & Leur \\
        \hline
    \end{tabular}
\end{center}
For \textit{Me, Te}, they are going to become \textit{Moi, Toi} when they are placed after the verb (in infinitve sentences). Can be seen as we don't usually have a [e] sound at the end of the verb.

\subsection{En and Y as Object Pronouns}
\begin{itemize}
    \item \textit{En} translates to \textit{some, any, of it, of them}, usually replaces a noun introduced by \textit{de}.
    \item \textit{Y} translates to \textit{there}, usually replaces a noun introduced by \textit{à}.
\end{itemize}

\section{Reflexive Verbs with L'Impératif}
Use present tense with \textbf{tu, nous, vous}, for example \textit{se dépêcher} (to hurry):
\begin{center}
    Tu te dépêches $\rightarrow$ Dépêche-toi \\
    Nous nous dépêchons $\rightarrow$ Dépêchons-nous \\
    Vous vous dépêchez $\rightarrow$ Dépêchez-vous
\end{center}

\subsection{Negation}
\begin{center}
    Ne te dépêche pas (Don't hurry) \\
    Ne nous dépêchons pas (Don't hurry) \\
    Ne vous dépêchez pas (Don't hurry)
\end{center}

\chapter{Infinitive}

\begin{enumerate}
    \item Verb + infinitive:
    \begin{center}
        J'aime \textcolor{red}{chanter} (I like to sing) \\
        J'aimais \textcolor{red}{chanter} (I enjoyed singing) \\
        J'aimerais \textcolor{red}{chanter} (I would like to sing)
    \end{center}
    Some phrases, only the first verb will be conjugated for the following list
    \begin{itemize}
        \item Aimer \underline{faire} qqch (to like to do something) 
        \item Devoir \underline{faire} qqch (to have to do something) 
        \item Oser \underline{faire} qqch (to dare to do something) 
        \item Pouvoir \underline{faire} qqch (to be able to do something) 
        \item Souhaiter \underline{faire} qqch (to wish to do something)
        \item Venir \underline{faire} qqch (to come to do something)
    \end{itemize}

    \item Verb + à + infinitive:
    \begin{itemize}
        \item Aspirer à \underline{faire} qqch (to aspire to do something)
        \item Continuer à \underline{faire} qqch (to continue to do something)
        \item S'engager à \underline{faire} qqch (to commit to do something)
        \item Persister à \underline{faire} qqch (to persist to do something)
        \item Inviter qqn à \underline{faire} qqch (to invite someone to do something)
        \item Parvenir à \underline{faire} qqch (to manage / succeed to do something)
    \end{itemize}

    \item Verb + à + noun:
    \begin{itemize}
        \item Aller à qqn (to suit someone)
        \item Croire à qqch (to believe in something)
        \item Être à qqn (to belong to someone)
        \item Se fier à qqn / qqch (to trust someone / something)
        \item Jouer à qqch (to play something)
        \item Obéir à qqn (to obey someone) 
    \end{itemize}

    \item Verb + de + infinitive:
    \begin{itemize}
        \item Accepter de \underline{faire} qqch (to accept to do something)
        \item Avoir envie de \underline{faire} qqch (to feel like doing something)
        \item Dire à qqn de \underline{faire} qqch (to tell someone to do something)
        \item Oublier de \underline{faire} qqch (to forget to do something)
        \item Refuser de \underline{faire} qqch (to refuse to do something)
        \item Rêver de \underline{faire} qqch (to dream of doing something)
    \end{itemize}

    \item Verb + de + noun:
    \begin{itemize}
        \item S'agir de qqn / qqch (to be about someone / something)
        \item Avoir envie de qqch (to want something)
        \item Se passer de qqn / qqch (to do without someone / something)
        \item Rêver de qqn / qqch (to dream of someone / something)
        \item Vivre de qqch (to live on something)
    \end{itemize}
\end{enumerate}
à and de are linked to the first verb, not the second.

\section{Where to Add Infinitive}

\begin{enumerate}
    \item Il est + adjective + de + infinitive (very formal way):
    \begin{itemize}
        \item Since Être is a non-personal verb, so the adjective will always stay masculine singular.
        \item Examples:
        \begin{itemize}
            \item Il est dangereux de (It is dangerous to)
            \begin{center}
                Il est dangereux de \textcolor{red}{jouer} avec le feu \\
                (It's dangerous to play with fire)  
            \end{center}
            \item Il est interdit de (It is forbidden to)
            \begin{center}
                Il est interdit de \textcolor{red}{fumer} à l'intérieur \\
                (It's forbidden to smoke inside)
            \end{center}
            \item Il est facile de (It is easy to)
            \item Il est impossible de (It is impossible to) 
            \begin{center}
                Il est impossible de \textcolor{red}{réparaer} cette voiture \\
                (It is impossible to repair this car)
            \end{center}
            \item Il est possible de (It is possible to)
            \item Il est inutile de (It is useless to)
        \end{itemize}
    \end{itemize}
    \item C'est  + adjective + de + infinitive (informal way, spoken French):
    \begin{itemize}
        \item The adjective will stay masculine singular.
        \item Examples:
        \begin{itemize}
            \item C'est interdit de (It is forbidden to)
            \item C'est difficile de (It is difficult to)
            \begin{center}
                C'est difficile de \textcolor{red}{se détendre} \\
                (It's difficult to relax)
            \end{center}
            \item C'est bien de (It is good to) 
            \begin{center}
                C'est bien de \textcolor{red}{changer} d'air \\
                (It's good to get some fresh air)
            \end{center}
            \item C'est dommage de (It's too bad to)
        \end{itemize}
    \end{itemize}
    \item C'est + adjective + à + infinitive:
    \begin{itemize}
        \item Use this construction to refer to something previously mentioned in the conversation.
        \item Examples:
        \begin{itemize}
            \item C'est bon à \textcolor{red}{savoir} (It's good to know)
            \item C'est difficile à \textcolor{red}{dire} (It's difficult to say)
            \item C'est impossible à \textcolor{red}{dire} (It's impossible to say)
        \end{itemize}
    \end{itemize}
    \item $\dots$ seul\_ + à + infinitive, examples:
    \begin{itemize}
        \item Le seul à (m) (the only one to)
        \begin{center}
            Il est le seul à \textcolor{red}{comprendre} son écriture \\
            (He is the only one to understand his writing)
        \end{center}
        \item La seule à (f) (the only one to)   
        \item Les seuls à (m) (the only ones to)
        \item Les seules à (f) (the only ones to)
        \begin{center}
            Les seules à \textcolor{red}{avoir} la réponse sont ses tantes \\
            (The only ones to have the answer are her aunts)
        \end{center}
        \item Nombreux à (m) (many are)
        \begin{center}
            Ils sont nombreux à \textcolor{red}{attendre} des nouvelles \\
            (Many are waiting for news)  
        \end{center}
        \item Nombreuses à (f) (many are)
        \begin{center}
            Elles sont nombreuses a \textcolor{red}{arrêter} leurs études \\
            (Many are stopping their studies)
        \end{center}
    \end{itemize}
    \item Ordinal number + à + infinitive
    \item Examples
    \begin{itemize}
        \item Le premier à (m) (the first one to)
        \begin{center} 
            Il est le premier à \textcolor{red}{aller} à l'université \\
            (He is the first one to go to university)
        \end{center}
        \item La première à (f) (the first one to)
        \item Les premiers à (m) (the first ones to)
        \item Les premières à (f) (the first ones to)
        \item Le deuxième à (m) (the second one to)
        \item La deuxième à (f) (the second one to)
        \item etc.
    \end{itemize}
    \item Il faut + infinitive:
    \begin{itemize}
        \item Il faut + que = subjonctif
        \item Il faut + verb = infinitive
        \begin{center}
            Il faut \textcolor{red}{partir} maintenant \\
            (We have to leave now) \\
            Il fault \textcolor{red}{aller} faire des courses \\
            (We have to go grocery shopping) 
        \end{center}
    \end{itemize}
    \item The past infinitive:
    \begin{itemize}
        \item Compound tense.
        \item Avoir + past participle
        \item Être + past participle
        \begin{center}
            \begin{tabular}{|c|c|}
                \hline
                \textbf{Infinitif présent} & \textbf{Infinitif passé} \\
                \hline
                Chanter & Avoir chanté \\
                Partir & Être parti \\
                Se préparer & S'être préparé \\
                \hline
            \end{tabular}
        \end{center}
        \item For conjugations with avoir and être, refer to the past participle section.
        \item The verbs always give the same auxiliary no matter the tense, for example:
            \begin{center}
                Nous vous remercions d'\textcolor{red}{être venus} (We thank you for coming) \\
            \end{center}
        \item After the preposition \textit{Après}, the infinitive is always a past infinitive, for example:
            \begin{center}
                Je suis partie après \textcolor{red}{avoir rendu} les clés \\
                (I left after returning the keys) \\
                Il est fatigué après \textcolor{red}{avoir mangé} \\
                (He is tired after eating)  
            \end{center}
        \item After the verbs \textit{Se rappeler de} or \textit{Se souvenir de} (to remember), the infinitive is always a past infinitive, since we are talking about something in the past, for example:
            \begin{center}
                Je ne me rappelle pas d'\textcolor{red}{avoir commandé} ça \\
                (I left after returning the keys) \\
                Est-ce que tu te souviens d'\textcolor{red}{avoir fermé} la porte? \\
                (Do you remember closing the door?)
            \end{center}
    \end{itemize}
    \item The negative infinitive: 
    \begin{itemize}
        \item Both \textit{ne} and \textit{pas} are placed before the infinitive, not matter if its a past or present infinitive.
        \begin{center}
            J'essaye de \underline{ne pas} \textcolor{red}{fumer} \\ 
            (I try not to smoke) \\
            C'est sale de \underline{ne pas} \textcolor{red}{se laver} les mains \\
            (It's dirty not to wash your hands)
        \end{center}
        \item If we have a pronoun, an direct object pronoun and an indirect object pronoun, the pronoun is going to stay between the negation and the verb:
        \begin{center}
            Je préfère ne pas \underline{y} \textcolor{red}{retourner} \\
            (I prefer to not go back (there)) \\
            C'est difficile de ne pas \underline{le} \textcolor{red}{voir} \\
            (It's difficult to not see it)
        \end{center}
    \end{itemize}
    \item Merci de + infinitive:
    \begin{itemize}
        \item Very formal way, usually find in signs.
        \begin{center}
            Merci d'\textcolor{red}{arriver} tôt \\
            (Thank you for arriving early) \\
            Merci de \textcolor{red}{fermer} les fenêtres \\
            (Thank you for closing the windows) \\
            Merci d'\textcolor{red}{étre venu} \\
            (Thank you for coming) 
        \end{center}
    \end{itemize}
\end{enumerate}

\chapter{Le Subjonctif}
Always after \textit{que}, used to express doubt, wishes, emotions, etc. Take the conjugation of plural \textit{ils} form, remove the -ent, and add the following endings:
    \begin{itemize}
        \item -e
        \item -es
        \item -e
        \item -ions
        \item -iez
        \item -ent
    \end{itemize}

\section{Conjugation}
    The verb groups:
        \begin{enumerate}
            \item Donner (to give): Ils donnent
                \begin{center}
                    \begin{tabular}{|c|}
                        \hline
                        Que je donn\textcolor{red}{e} \\
                        \hline  
                        Que tu donn\textcolor{red}{es} \\
                        \hline
                        Qu'il/elle/on donn\textcolor{red}{e} \\
                        \hline
                        Que nous donn\textcolor{red}{ions} \\
                        \hline
                        Que vous donn\textcolor{red}{iez} \\
                        \hline
                        Qu'ils/elles donn\textcolor{red}{ent} \\
                        \hline
                    \end{tabular}
                \end{center}

            \item Finir (to finish): Ils finissent
                \begin{center}
                    \begin{tabular}{|c|}
                        \hline
                        Que je finiss\textcolor{red}{e} \\
                        \hline  
                        Que tu finiss\textcolor{red}{es} \\
                        \hline
                        Qu'il/elle/on finiss\textcolor{red}{e} \\
                        \hline
                        Que nous finiss\textcolor{red}{ions} \\
                        \hline
                        Que vous finiss\textcolor{red}{iez} \\
                        \hline
                        Qu'ils/elles finiss\textcolor{red}{ent} \\
                        \hline
                    \end{tabular}
                \end{center}

            \item Attendre (to wait): Ils attendent
                \begin{center}
                    \begin{tabular}{|c|}
                        \hline
                        Que j'attend\textcolor{red}{e} \\
                        \hline
                        Que tu attend\textcolor{red}{es} \\
                        \hline
                        Qu'il/elle/on attend\textcolor{red}{e} \\
                        \hline
                        Que nous attend\textcolor{red}{ions} \\
                        \hline
                        Que vous attend\textcolor{red}{iez} \\
                        \hline
                        Qu'ils/elles attend\textcolor{red}{ent} \\
                        \hline                
                    \end{tabular}
                \end{center}

            \item Verbs with spelling changes in the stem, the ones changes with \textit{nous, vous}, the spelling change in the stem stays:
                \begin{center}
                    \begin{tabular}{|c|c|c|}
                        \hline
                        \textbf{Essuyer} Ils essuient & \textbf{Peser} Ils pèsent & \textbf{Célébrer} Ils célèbrent \\
                        \hline
                        Que j'essui\textcolor{red}{e} & Que je pès\textcolor{red}{e} & Que je célèbr\textcolor{red}{e} \\
                        \hline
                        Que tu essui\textcolor{red}{es} & Que tu pès\textcolor{red}{es} & Que tu célèbr\textcolor{red}{es} \\
                        \hline
                        Qu'il/elle/on essui\textcolor{red}{e} & Qu'il/elle/on pès\textcolor{red}{e} & Qu'il/elle/on célèbr\textcolor{red}{e} \\
                        \hline
                        Que nous essu\textcolor{red}{yions} & Que nous p\textcolor{red}{e}s\textcolor{red}{ions} & Que nous cél\textcolor{red}{é}br\textcolor{red}{ions} \\
                        \hline
                        Que vous essu\textcolor{red}{yiez} & Que vous p\textcolor{red}{e}s\textcolor{red}{iez} & Que vous cél\textcolor{red}{é}br\textcolor{red}{iez} \\
                        \hline
                        Qu'ils/elles essu\textcolor{red}{ent} & Qu'ils/elles pes\textcolor{red}{ent} & Qu'ils/elles célébr\textcolor{red}{ent} \\
                        \hline
                    \end{tabular} 
                \end{center}

            \item Irregular verbs with a regular sobjonctif (the following examples works for all verbs conjugated the same way):
                \begin{itemize}
                    \item -TTRE: Battre (to beat)
                    \begin{center}
                        Ils battent \\
                        \begin{tabular}{|c|}
                            \hline
                            Que je batt\textcolor{red}{e} \\
                            \hline
                            Que nous batt\textcolor{red}{ions} \\
                            \hline
                            Qu'ils batt\textcolor{red}{ent} \\
                            \hline
                        \end{tabular}
                    \end{center}

                    \item Takes an 's' in the present tense stem: Conduire (to drive)
                    \begin{center}
                        Ils conduisent \\
                        \begin{tabular}{|c|}
                            \hline
                            Que je conduis\textcolor{red}{e} \\
                            \hline
                            Que nous conduis\textcolor{red}{ions} \\
                            \hline
                            Qu'ils conduis\textcolor{red}{ent} \\
                            \hline
                        \end{tabular}
                    \end{center}

                    \item Takes an 'v' in the stem: Écrire (to write)
                    \begin{center}
                        Ils écrivent \\
                        \begin{tabular}{|c|}
                            \hline
                            Que j'écriv\textcolor{red}{e} \\
                            \hline
                            Que nous écriv\textcolor{red}{ions} \\
                            \hline
                            Qu'ils écriv\textcolor{red}{ent} \\
                            \hline
                        \end{tabular}
                    \end{center}
                \end{itemize}
            
            \item Irregular verbs with 2 different stems:
                \begin{center}
                    \begin{tabular}{|c|c|c|c|}
                        \hline
                        & \textit{je, tu, il/elle/on, ils/elles} & \textit{nous, vous} & \textbf{English} \\
                        \hline
                        Boire (Ils boivent) & Que je boive & Que nous b\textcolor{red}{u}vions & To drink \\
                        \hline
                        Devoir (Ils doivent) & Que je doive & Que nous d\textcolor{red}{e}vions & To have to \\
                        \hline
                        Envoyer (Ils envoient) & Que j'envoie & Que nous envo\textcolor{red}{y}ions & To send \\
                        \hline
                        Mourir (Ils meurent) & Que je meure & Que nous m\textcolor{red}{ourr}ions & To die \\
                        \hline
                        Recevoir (Ils reçoivent) & Que je reçoive & Que nous re\textcolor{red}{ce}vions & To receive \\
                        \hline
                        Tenir (Ils tiennent) & Que je tienne & Que nous t\textcolor{red}{en}ions & To hold \\
                        \hline
                        Venir (Ils viennent) & Que je vienne & Que nous v\textcolor{red}{en}ions & To come \\
                        \hline
                        Voir (Ils voient) & Que je voie & Que nous vo\textcolor{red}{y}ions & To see \\
                        \hline
                    \end{tabular}
                \end{center}
            
            \item Irregular with an irregular stem:
                \begin{itemize}
                    \item Faire (to do / to make):
                        \begin{center}
                            \begin{tabular}{|c|}
                                \hline
                                Que je fass\textcolor{red}{e} \\
                                \hline
                                Que tu fass\textcolor{red}{es} \\
                                \hline
                                Qu'il/elle/on fass\textcolor{red}{e} \\
                                \hline
                                Que nous fass\textcolor{red}{ions} \\
                                \hline
                                Que vous fass\textcolor{red}{iez} \\
                                \hline
                                Qu'ils/elles fass\textcolor{red}{ent} \\
                                \hline
                            \end{tabular}
                        \end{center}

                    \item Savoir (to know):
                        \begin{center}
                            \begin{tabular}{|c|}
                                \hline
                                Que je sach\textcolor{red}{e} \\
                                \hline
                                Que tu sach\textcolor{red}{es} \\
                                \hline
                                Qu'il/elle/on sach\textcolor{red}{e} \\
                                \hline
                                Que nous sach\textcolor{red}{ions} \\
                                \hline
                                Que vous sach\textcolor{red}{iez} \\
                                \hline
                                Qu'ils/elles sach\textcolor{red}{ent} \\
                                \hline
                            \end{tabular}
                        \end{center}

                    \item Pouvoir (to be able to):
                        \begin{center}
                            \begin{tabular}{|c|}
                                \hline
                                Que je puiss\textcolor{red}{e} \\
                                \hline
                                Que tu puiss\textcolor{red}{es} \\
                                \hline
                                Qu'il/elle/on puiss\textcolor{red}{e} \\
                                \hline
                                Que nous puiss\textcolor{red}{ions} \\
                                \hline
                                Que vous puiss\textcolor{red}{iez} \\
                                \hline
                                Qu'ils/elles puiss\textcolor{red}{ent} \\
                                \hline
                            \end{tabular}
                        \end{center}

                    \item Aller (to go): -aill for everything except for \textit{nous, vous}, -all for \textit{nous, vous}
                        \begin{center}
                            \begin{tabular}{|c|}
                                \hline
                                Que j'ail\textcolor{red}{le} \\
                                \hline
                                Que tu aill\textcolor{red}{es} \\
                                \hline
                                Qu'il/elle/on aill\textcolor{red}{e} \\
                                \hline
                                Que nous all\textcolor{red}{ions} \\
                                \hline
                                Que vous all\textcolor{red}{iez} \\
                                \hline
                                Qu'ils/elles aill\textcolor{red}{ent} \\
                                \hline
                            \end{tabular}
                        \end{center}

                    \item Valoir (to be worth): -vaill for everything except for \textit{nous, vous}, -val for \textit{nous, vous}
                        \begin{center}
                            \begin{tabular}{|c|}
                                \hline
                                Que je vaill\textcolor{red}{e} \\
                                \hline
                                Que tu vaill\textcolor{red}{es} \\
                                \hline
                                Qu'il/elle/on vaill\textcolor{red}{e} \\
                                \hline
                                Que nous val\textcolor{red}{ions} \\
                                \hline
                                Que vous val\textcolor{red}{iez} \\
                                \hline
                                Qu'ils/elles vaill\textcolor{red}{ent} \\
                                \hline
                            \end{tabular}
                        \end{center}

                    \item Vouloir (to want): -euill for everything except for \textit{nous, vous}, -ouil for \textit{nous, vous}
                        \begin{center}
                            \begin{tabular}{|c|}
                                \hline
                                Que je veuill\textcolor{red}{e} \\
                                \hline
                                Que tu veuill\textcolor{red}{es} \\
                                \hline
                                Qu'il/elle/on veuill\textcolor{red}{e} \\
                                \hline
                                Que nous voul\textcolor{red}{ions} \\
                                \hline
                                Que vous voul\textcolor{red}{iez} \\
                                \hline
                                Qu'ils/elles veuill\textcolor{red}{ent} \\
                                \hline
                            \end{tabular}
                        \end{center}

                    \item Être and Avoir (will be used for subjonctif passé):
                        \begin{center}
                            \begin{tabular}{|c|c|}
                                \hline
                                \textit{Être} & \textit{Avoir} \\
                                \hline
                                Que je sois & Que j'aie \\
                                \hline
                                Que tu sois & Que tu aies \\
                                \hline
                                Qu'il/elle/on soit & Qu'il/elle/on ait \\
                                \hline
                                Que nous soyons & Que nous ayons \\
                                \hline
                                Que vous soyez & Que vous ayez \\
                                \hline
                                Qu'ils/elles soient & Qu'ils/elles aient \\
                                \hline
                            \end{tabular}
                        \end{center}
                \end{itemize}
        \end{enumerate}

    \subsection{Le Subjonctif Passé}
        The structure is:
        \begin{center}
            Être / Avoir + past participle
        \end{center}
        \begin{itemize}
            \item Verbs conjugated with \textit{être} or \textit{avoir} stays the same, for example:
                \begin{enumerate}
                    \item Verb with \textit{avoir}:
                        \begin{center}
                            Demander (to ask) PP: Demandé \\
                            \begin{tabular}{|c|}
                                \hline
                                Que j'aie demandé \\
                                \hline
                                Que tu aies demandé \\
                                \hline
                                Qu'il/elle/on ait demandé \\
                                \hline
                                Que nous ayons demandé \\
                                \hline
                                Que vous ayez demandé \\
                                \hline
                                Qu'ils/elles aient demandé \\
                                \hline
                            \end{tabular}
                        \end{center}

                    \item Verb with \textit{être}:
                        \begin{center}
                            Partir (to leave) PP: Parti \\
                            \begin{tabular}{|c|}
                                \hline
                                Que je sois parti(e) \\
                                \hline
                                Que tu sois parti(e) \\
                                \hline
                                Qu'il soit parti \\
                                \hline
                                Qu'elle soit partie \\
                                \hline
                                Qu'on soit parti(e)s \\
                                \hline
                                Que nous soyons parti(e)s \\
                                \hline
                                Que vous soyez parti(e)s \\
                                \hline
                                Qu'ils soient partis \\
                                \hline
                                Qu'elles soient parties \\
                                \hline
                            \end{tabular}
                        \end{center}
                        Note add (e) if subject is feminine.
                \end{enumerate}

            \item Reflexive verb are conjugated with \textit{être}, for example:
                \begin{center}
                    Se trouver (to be) as a location. PP: Trouvé \\
                    \begin{tabular}{|c|}
                        \hline
                        Que je me sois trouvé(e) \\
                        \hline
                        Que tu te sois trouvé(e) \\
                        \hline
                        Qu'il se soit trouvé \\
                        \hline
                        Qu'elle se soit trouvée \\
                        \hline
                        Qu'on se soit trouvé(e)s \\
                        \hline
                        Que nous nous soyons trouvé(e)s \\
                        \hline
                        Que vous vous soyez trouvé(e)s \\
                        \hline
                        Qu'ils se soient trouvés \\
                        \hline
                        Qu'elles se soitent trouvées \\
                        \hline
                    \end{tabular}
                \end{center}
        \end{itemize}

\section{Quand Utiliser le Subjonctif}
    The subjonctif is used after \textit{que}, mostly in subordinate clause and indicates a wish, a regret, an emotion, an opinion, or a doubt. Subordinate means that the clause cannot stand alone as a sentence. For example,
        \begin{center}
            Je veux que ce problème \textcolor{red}{soit réglé} \\
            (I want this problem to be solved)
        \end{center}
        \begin{itemize}
            \item \textit{Je veux} is the main clause
            \item \textit{Ce problème soit réglé} is the subordinate clause
            \item \textit{soit réglé} is the subjonctif
            \item \textit{que} is not always translated to \textit{that} in English
        \end{itemize}
        
    \subsection{WEIRDO}
    To memorize scenarios where the subjonctif is used, remember the acronym 
        \begin{center}
            WEIRDO
        \end{center}
        \begin{itemize}
            \item W: Wish, want, will
                \begin{center}
                    \begin{tabular}{|c|c|}
                        \hline
                        \textbf{French} & \textbf{English} \\
                        \hline
                        Attendre que & To wait for \\
                        \hline
                        Avoir envie que & To want that \\
                        \hline
                        Insister pour que & To insist that \\
                        \hline
                        Demander que & To ask that \\
                        \hline
                    \end{tabular}
                \end{center}
                \begin{center}
                    J'insiste pour qu'elle \textcolor{red}{reçoive} un cadeau aussi \\
                    (I insist that she gets a gift as well)
                \end{center}
                It's not always subjonctif if we have \textit{que}, consider the following example (this is a indicative, not subjonctif):
                \begin{center}
                    J'espère que tu seras là (I hope you will be there)
                \end{center}
                
            \item E: Emotion
                \begin{itemize}
                    \item Takes the form:
                    \begin{center}
                        Être + adjective + que
                    \end{center}

                    \item adjective agrees in gender.
                    \begin{center}
                        \begin{tabular}{|c|c|}
                            \hline
                            \textbf{French} & \textbf{English} \\
                            \hline
                            Être content / contente que & To be happy that \\
                            \hline
                            Être étonné / étonnée que & To be surprised that \\
                            \hline
                            Être soulagé / soulagée que & To be relieved that \\
                            \hline
                            Être desolé / desolée que & To be sorry that \\
                            \hline
                        \end{tabular}
                    \end{center}

                    \item Examples:
                    \begin{enumerate}
                        \item Subjonctif Passé:
                        \begin{center}
                            Il était tellement content que tu \textcolor{red}{sois venue} \\
                            (He was so happy that you came) 
                        \end{center}

                        \item Subjonctif Présent:
                        \begin{center}
                            Je suis ravie que tu \textcolor{red}{viennes} pour le week-end \\
                            (I'm delighted that you're coming for the weekend)
                        \end{center}
                    \end{enumerate}

                    \item Subjonctif is not used if the adjective is followed by a infinitive verb, for example:
                    \begin{center}
                        Nous sommes désolés d'apprendre la nouvelle \\
                        (We are sorry to hear the news)
                    \end{center}
                    \begin{itemize}
                        \item There is no \textit{que} in the sentence, so it's not subjonctif
                    \end{itemize}
                \end{itemize}

            \item I: Impersonal expressions used to talk about obligations
            \begin{itemize}
                \item Most will be with (100\% subjonctif if present):
                \begin{center}
                    Il faut que (It is necessary that / to have to)
                \end{center}

                \item If there is no \textit{que}, \textit{il faut} plus the verb will be the infinitive.
                \begin{center}
                    \begin{tabular}{|c|c|}
                        \hline
                        \textbf{French} & \textbf{English} \\
                        \hline
                        Il faut que & It is necessary that \\
                        \hline
                        Il est important que & It is important that \\
                        \hline
                        Il est préférable que & It is preferable that \\
                        \hline
                        Il vaut mieux que & It is better that \\
                        \hline
                    \end{tabular}
                \end{center}

                \item Two examples with \textit{il faut}:
                \begin{enumerate}
                    \item Il faut que j'\textcolor{red}{achète} à manger (I have to buy food)
                    \item Il faut que tu \textcolor{red}{prennes} tes médicaments (You have to take your medicine)
                \end{enumerate}

                \item Example with \textit{il vaut mieux}, here is (S'en aller):
                \begin{center}
                    Il vaut mieux que tu \textcolor{red}{t'en ailles} (It's better that you leave)
                \end{center}
            \end{itemize}
            \item R: Request, recommendation, requirement
            \begin{center}
                \begin{tabular}{|c|c|}
                    \hline
                    \textbf{French} & \textbf{English} \\
                    \hline
                    Recommander que & To recommend that \\
                    \hline
                    Regretter que & To regret that \\
                    \hline
                    Suggérer que & To suggest that \\
                    \hline
                \end{tabular} \\
                \vspace{0.2cm}
                Le ministre recommande que les vaccins \textcolor{red}{soient distribués} au plus vite \\
                (The minister recommends that vaccines be distributed as soon as possible) \\
                Je regrette que tu l'\textcolor{red}{apprennes} de cette façon \\
                (I regret that you learn it this way)
            \end{center}
            \item D: Doubt, denial, disbelief
            \begin{center}
                \begin{tabular}{|c|c|}
                    \hline
                    \textbf{French} & \textbf{English} \\
                    \hline
                    Il est possible que & It is possible that \\
                    \hline
                    Il est incroyable que & It is unbelievable that \\
                    Il semble que & It seems that \\
                    \hline
                \end{tabular}
            \end{center}
            \begin{itemize}
                \item \textit{Il est} can be replaced by \textit{c'est}
                \item \textit{Il est} is more in books, formal french, where \textit{c'est} is more spoken french, informal.
            \end{itemize}
            \begin{center}
                Il semble que la tempête \textcolor{red}{soit poassée} \\
                (It seems that the storm has passed) \\
                C'est impossible que tu \textcolor{red}{aies raté} ça! \\
                (It's impossible that you missed that!)
            \end{center}
            Note on \textit{Croire, Avoir l'impression, Penser}:
            \begin{itemize}
                \item These verbs can also express a doubt, all of them used in the positive (in the informative sentence) are going to be used with the \textbf{indicative}.
                \item Only used in subjonctif in negative or interrogative sentences (inversion question).
                \item Examples for each:
                \begin{enumerate}
                    \item Croire 
                    \begin{center}
                        \underline{Indicatif} \\
                        Je crois qu'il \textcolor{red}{sait} cuisiner \\
                        (I think he knows how to cook) \\
                        \underline{Subjonctif} \\
                        Je ne crois pas qu'il \textcolor{red}{sache} cuisiner \\
                        (I don't think he knows how to cook)
                    \end{center}

                    \item Avoir l'impression
                    \begin{center}
                        \underline{Indicatif} \\
                        J'ai l'impression qu'elle \textcolor{red}{est} déçue \\
                        (I feel like she is disappointed) \\
                        \underline{Subjonctif} \\
                        Je n'ai pas l'impression qu'elle \textcolor{red}{soit} déçue \\
                        (I don't feel like she is disappointed)
                    \end{center}

                    \item Penser
                    \begin{center}
                        \underline{Indicatif} \\
                        Je pense qu'il \textcolor{red}{a} du talent \\
                        (I think he has talent) \\
                        \underline{Subjonctif} \\
                        Je ne pense pas qu'il \textcolor{red}{ait} du talent \\
                        (I don't think he has talent)
                    \end{center}
                \end{enumerate}
                
            \end{itemize}
            
            \item O: Opinions
            \begin{center}
                \begin{tabular}{|c|c|}
                    \hline
                    \textbf{French} & \textbf{English} \\
                    \hline
                    Il est bien que & It is good that \\
                    \hline
                    Il est dommage que & It is a shame that \\
                    \hline
                    Il est logique que & It is logical that \\
                    \hline
                    Il est rare que & It is rare that \\
                    \hline
                \end{tabular}
            \end{center}
            \begin{itemize}
                \item \textit{Il est} can be replaced by \textit{c'est} in all cases once again.
                \item The gender of the adjective does not change after \textit{c'est} or \textit{il est}, always masculine singular.
            \end{itemize}
            \begin{center}
                C'est dommage que vous ne \textcolor{red}{puissiez} pas venir \\
                (It's a shame that you can't come) \\
                C'est logique qu'il \textcolor{red}{ait} gagné \\
                (It's logical that he won)
            \end{center}
        \end{itemize}

    \subsection{After some Conjunctions}
        \begin{center}
            \begin{tabular}{|c|c|}
                \hline
                \textbf{French} & \textbf{English} \\
                \hline
                À condition que & Provided that / if \\
                \hline
                Afin que & So that \\
                \hline
                Jusqu'à ce que & Until \\
                \hline
                Pour que & So that \\
                \hline
            \end{tabular}
        \end{center}
        Example:
        \begin{center}
            Je suis d'accord à condition que tu \textcolor{red}{te tiennes} bien \\
            (I agree if you behave) \\
            Elle utilise son vieux téléphone jusqu'à ce qu'elle \textcolor{red}{reçoive} le nouveau \\
            (She uses her old phone until she receives the new one)
        \end{center}

    \subsection{The "Ne" Explétif}
        \begin{itemize}
            \item Doesn't add any meaning
            \item Doesn't turn the sentence into a negative sentence
            \item Only designed to emphasize the first part of the sentence
            \item Example phrases:
                \begin{center}
                    \begin{tabular}{|c|c|}
                        \hline
                        \textbf{French} & \textbf{English} \\
                        \hline
                        À moins que ... (ne) & Unless \\
                        \hline
                        Avant que ... (ne) & Before that \\
                        \hline
                        De crainte que ... (ne) & For fear that \\
                        \hline
                        De peur que ... (ne) & For fear that \\
                        \hline
                        Sans que ... (ne) & Without \\
                        \hline
                    \end{tabular}
                \end{center}
            \item Example sentences:
                \begin{center}
                    Il marche tout seul sans que je \textcolor{red}{ne} l'\textcolor{red}{aide} \\
                    (He walks alone without me helping him) \\
                    Prépare le dossier avant que je \textcolor{red}{n'arrive} \\
                    (Get the file ready before I arrive)
                \end{center}
        \end{itemize}

    \subsection{After the superlative}
        \begin{enumerate}
            \item Le seul / La seule ... qui / que / dont (The only one that)
                \begin{center}
                    C'est le seul qui \textcolor{red}{sache} faire fonctionner cette machine \\
                    (He is the only one who knows how to operate this machine)
                \end{center}

            \item Le premier / La première ... qui / que / dont (The first one that)
                \begin{center}
                    La première qui \textcolor{red}{arrive} gagne le gros lot \\
                    (The first one who arrives wins the big prize)
                \end{center}

            \item Le meilleur / La meilleure ... qui / que (The best (one) that)
                \begin{center}
                    C'est la meilleure tarte que j'\textcolor{red}{aie} jamais \textcolor{red}{mangée} \\
                    (It's the best pie I've ever eaten)
                \end{center}
            
            \item Il n'y a que ... qui / que (There is only ... that)
                \begin{center}
                    Il n'y a que l'argent qui \textcolor{red}{compte} pour toi \\
                    (There is only money for you)
                \end{center}
            
            \item Il y a peu de ... qui / que (There aren't many ... that)
                \begin{center}
                    Il y a peu d'oisillons qui \textcolor{red}{survivent} le premier mois \\
                    (There aren't many chicks that survive the first month)
                \end{center}
        \end{enumerate}

\chapter{Le Plus-Que-Parfait}
    \begin{itemize}
        \item Same as \textit{I had done} in English
        \item The structure is:
            \begin{center}
                Avoir ou Être conjugated in the imparfait + past participle
            \end{center}
        \item Used to talk about an action or a situation that happened \underline{before} another past action.
        \item If it's alone, it means we already know the action that happened in the past.
    \end{itemize}

    \section{Conjugation of Avoir and Être in Imparfait}
        \begin{center}
            \begin{tabular}{|c|c|}
                \hline
                \textbf{Avoir} & \textbf{Être} \\
                \hline
                J'avais & J'étais \\
                \hline
                Tu avais & Tu étais \\
                \hline
                Il/elle/on avait & Il/elle/on était \\
                \hline
                Nous avions & Nous étions \\
                \hline
                Vous aviez & Vous étiez \\
                \hline
                Ils/elles avaient & Ils/elles étaient \\
                \hline
            \end{tabular}
        \end{center}
    
    \section{Common Verbs Conjugation}
        \begin{enumerate}
            \item Téléphoner (to call):
                \begin{center}
                    \begin{tabular}{|c|}
                        \hline
                        J'avais téléphoné \\
                        \hline
                        Tu avais téléphoné \\
                        \hline
                        Il/elle/on avait téléphoné \\
                        \hline
                        Nous avions téléphoné \\
                        \hline
                        Vous aviez téléphoné \\
                        \hline
                        Ils/elles avaient téléphoné \\
                        \hline
                    \end{tabular}
                \end{center}

            \item Partir (to leave), always conjugated with \textit{Être}:
                \begin{center}
                    \begin{tabular}{|c|}
                        \hline
                        J'étais parti(e) \\
                        \hline
                        Tu étais parti(e) \\
                        \hline
                        Il était parti \\
                        \hline
                        Elle était partie \\
                        \hline
                        On était parti(e)s \\
                        \hline
                        Nous étions parti(e)s \\
                        \hline
                        Vous étiez parti(e)s \\
                        \hline
                        Ils étaient partis \\
                        \hline
                        Elles étaient parties \\
                        \hline
                    \end{tabular}
                \end{center}
                Note the gender and number agreement applies for verb conjugated with \textit{Être}.
        \end{enumerate}

    \section{Reflexive Verbs}
        Always conjugated with \textit{Être}. For example:
        \begin{center}
            Se changer (to get changed)
            \begin{tabular}{|c|}
                \hline
                Je m'étais changé(e) \\
                \hline
                Tu t'étais changé(e) \\
                \hline
                Il s'était changé \\
                \hline
                Elle s'était changée \\
                \hline
                On s'était changé(e)s \\
                \hline
                Nous nous étions changé(e)s \\
                \hline
                Vous vous étiez changé(e)s \\
                \hline
                Ils s'étaient changés \\
                \hline
                Elles s'étaient changées \\
                \hline
            \end{tabular}
        \end{center}

    \section{Examples}
        \begin{center}
            J'\textcolor{red}{avais compris} quand le téléphone a sonné \\
            (I had understood when the phone rang) \\
            Vous \textcolor{red}{aviez} déjà \textcolor{red}{acheté} votre nouvelle voiture à ce moment-là \\
            (You had already bought your new car at that time) \\
            Je pensais que tu n'\textcolor{red}{avais} jamais \textcolor{red}{vu} ce film \\
            (I thought you had never seen this movie)
        \end{center}

\chapter{Futur Antérieur}
    \begin{itemize}
        \item Compound tense, so there are two differnet parts.
        \item Equivalent to \textit{I will have done} in English.
        \item The structure is
            \begin{center}
                Avoir ou Être conjugated in the futur simple + past participle
            \end{center}
        \item Used to talk about an action or a situation going to be completed \underline{before} another action in the future.
    \end{itemize}

    \section{Examples}
        \begin{center}
            Je \textcolor{red}{serai rentrée} quand tu \underline{arriveras} \\
            (I will have come home when you will arrive) \\
            Elle \underline{vivra} à l'étranger quand elle \textcolor{red}{aura fini} ses études \\
            (She will live abroad when she will be done with her studies) \\
            Quand il \underline{verra} l'oiseau, il \textcolor{red}{se sera envolé} \\
            (When he will see the bird, it will have flown away)
        \end{center}

    \section{Avoir and Être Conjugated in Futur Simple}
        \begin{center}
            \begin{tabular}{|c|c|}
                \hline
                \textbf{Avoir} & \textbf{Être} \\
                \hline
                J'aurai & Je serai \\
                \hline
                Tu auras & Tu seras \\
                \hline
                Il/elle/on aura & Il/elle/on sera \\
                \hline
                Nous aurons & Nous serons \\
                \hline
                Vous aurez & Vous serez \\
                \hline
                Ils/elles auront & Ils/elles seront \\
                \hline
            \end{tabular}
        \end{center}

\chapter{Conditionnel Passé}
    \begin{itemize}
        \item Equivalent to \textit{I would have done} in English.
        \item \textit{Would, should} are not used in French, instead it's conjugation of the auxiliary verb \textit{avoir} or \textit{être}.
        \item The structure is:
            \begin{center}
                Avoir ou Être conjugated in the conditionnel présent + past participle
            \end{center}
        \item Used to talk about a past situation, unsure information, a wish, or a regret.
        \item Very formal, often found in newspaper and news.
    \end{itemize}

    \section{Examples}
        \begin{center}
            D'après ce que l'on sait, les parents \textcolor{red}{auraient abandonné} l'enfant dans l'appartement. \\
            (From what we know, the parents would have left the child in the appartement) \\
            Il \textcolor{red}{serait mis} 10 jours pour arriver \\
            (It would have taken 10 days to arrive) \\
            Elle \textcolor{red}{aurait demandé} plusieurs fois un médicament \\
            (She would have asked for the medicine several times)
        \end{center}

    \section{Avoir and Être in Conditionnel Passé}
        \begin{center}
            \begin{tabular}{|c|c|}
                \hline
                \textbf{Avoir} & \textbf{Être} \\
                \hline
                J'aurais & Je serais \\
                \hline
                Tu aurais & Tu serais \\
                \hline
                Il/elle/on aurait & Il/elle/on serait \\
                \hline
                Nous aurions & Nous serions \\
                \hline  
                Vous auriez & Vous seriez \\
                \hline
                Ils/elles auraient & Ils/elles seraient \\
                \hline
            \end{tabular}
        \end{center}
  
\chapter{Le Participe Présent \& Le Gérondif}
    \begin{itemize}
        \item In French, it is very formal, for instance, in newspapers, formal speaking, news, books etc.
        \item We have 
            \begin{enumerate}
                \item Participe présent (rare in spoken French)
                \item Gérondif: EN + participe présent
            \end{enumerate}
    \end{itemize}

    \section{Le Participe Présent}
        \begin{itemize}
            \item Take stem from the present tense with \textit{Nous} and remove -ons and add -ant.
            \item There is no pronoun in front of the verb.
            \item For example:
                \begin{center}
                    \begin{tabular}{|c|c|c|}
                        \hline
                        \textbf{Present Tense Nous} & \textbf{Participe Présent} & \textbf{English} \\
                        \hline
                        Nous donn\textcolor{red}{\sout{ons}} & Donnant & Giving \\
                        \hline
                        Nous finiss\textcolor{red}{\sout{ons}} & Finissant & Finishing \\
                        \hline
                        Nous attend\textcolor{red}{\sout{ons}} & Attendant & Waiting \\
                        \hline
                        Nous prononç\textcolor{red}{\sout{ons}} & Prononçant & Pronouncing \\
                        \hline
                        Nous mange\textcolor{red}{\sout{ons}} & Mangeant & Eating \\
                        \hline
                        Nous craign\textcolor{red}{\sout{ons}} & Craignant & Fearing \\
                        \hline
                        \dots & & \\
                        \hline
                        Avoir & Ayant & Having \\
                        \hline
                        Être & Étant & Being \\
                        \hline
                        Savoir & Sachant & Knowing \\
                        \hline
                    \end{tabular}
                \end{center}
        \end{itemize}

        \subsection{Examples}
            The participe présent is referring to a noun. For example:
            \begin{center}
                \textcolor{red}{Voyant} Karen, il a changé de direction \\
                (Seeing Karen, he changed direction) \\
                \textcolor{red}{Habitant} près de la gare, il est constamment dérangé par le bruit des trains \\
                (Living by the train station, he is onstantlly bothered by the noise of trains)
            \end{center}

    \section{Le Gérondif}
        \begin{itemize}
            \item EN + Participe Présent
            \item Translates to \textit{while} or \textit{upon} in English.
            \item Not often used.
            \item To talk about an action related to the verb. It's usually about \underline{simultaneous} actions
                \begin{center}
                    J'écoute toujours de la musique \textcolor{red}{en travaillant} \\
                    (I always listen to music while working) \\
                    Nous avons rencontré Paul \textcolor{red}{en mangeant} au restaurant \\ 
                    (We met Paul while eating at the restaurant)
                \end{center}
            \item To talk about the cause and it mostly translates to \textit{by}
                \begin{center}
                    Elle a gagné 3 millions \textcolor{red}{en jouant} à la loterie \\
                    (She won 3 million by playing the lottery) \\
                    Il s'est fait mal \textcolor{red}{en tombant} de la chaise \\
                    (He hurt himself by falling off the chair)
                \end{center}
            \item Use the Gérondif without \textit{en} to replace a relative clause, usually includes \textit{qui}
                \begin{center}
                    Nous cherchons des guides \textcolor{red}{qui parlent} français \\
                    $\rightarrow$ Nous cherchons des guides \textcolor{red}{parlant} français \\
                    (We are looking for guides who speak French) \\
                    Les passagers \textcolor{red}{qui attendent} le train de 16 heures sont priés de changer de quai \\
                    $\rightarrow$ Les passagers \textcolor{red}{attendant} le train de 16 heures sont priés de changer de quai \\
                    (Passengers waiting for the 4 p.m. train are requested to change platforms)
                \end{center}
        \end{itemize}

    \section{Difference between the two Example}
        \begin{enumerate}
            \item Participe présent:
                \begin{center}
                    J'ai vu Nicole \textcolor{red}{traversant} la rue \\
                    (I saw Nicole crossing the street)
                \end{center}

            \item Gérondif, describes the action of self:
                \begin{center}
                    J'ai vu Nicole \textcolor{red}{en traversant} la rue \\
                    (I saw Nicole while crossing the street)
                \end{center}
        \end{enumerate}

\chapter{La Voix Passive}
    \begin{itemize}
        \item Used to shift the focus of the sentence to the second character or the second action, same as in English.
        \item Always conjugated with \textit{Être}, so the past participle will agree in gender and number with the subject.
        \item Not used as much in French, since it's very formal. French use \textit{On / Ils} in conversation, for example:
            \begin{center}
                La voiture a été réparée (The car has been fixed)
            \end{center}
            most likely to say 
            \begin{center}
                On a réparé la voiture (We fixed the car) \\
                Ils ont réparé la voiture (They fixed the car)
            \end{center}
            Note hear \textit{On} doesn't necessarily mean people we are speaking, it is referring to someone else who did it. So be careful when hearing \textit{On} in spoken French.
        \item One sentence in four different tense, goes from active to passive:
            \begin{enumerate}
                \item Présent:
                    \begin{center}
                        \textit{Actif.} Le banquier appelle la cliente (The banker is calling the client) \\
                        \textit{Passif.} La cliente \textcolor{red}{est appelée} par le banquier (The client is called by the banker)
                    \end{center}
                    Note we add \textit{par}, same as in English \textit{by}.

                \item Passé Composé:
                    \begin{center}
                        \textit{Actif.} Le banquier a appelé la cliente (The banker called the client) \\
                        \textit{Passif.} La cliente \textcolor{red}{a été appelée} par le banquier (The client has been called by the banker)
                    \end{center}

                \item Imparfait:
                    \begin{center}
                        \textit{Actif.}Le banquier appelait la cliente (The banker was calling the client) \\
                        \textit{Passif.}La cliente \textcolor{red}{était appelée} par le banquier (The client was called by the banker)
                    \end{center}

                \item Futur Simple:
                    \begin{center}
                        \textit{Actif.}Le banquier appellera la cliente (The banker will call the client) \\
                        \textit{Passif.}La cliente \textcolor{red}{sera appelée} par le banquier (The client will be called by the banker)
                    \end{center}
            \end{enumerate}
    \end{itemize}

\chapter{Le passé simple}
A tense that is not used anymore. Likely to find it in old books, poems, and formal writing. 

\end{document}